% Chapter 3, Section 2 _Linear Algebra_ Jim Hefferon
%  http://joshua.smcvt.edu/linearalgebra
%  2001-Jun-11
\section{Homomorfisme}
Definisi isomorfisme memiliki dua syarat.
Dalam bagian ini kita akan membahas syarat yang kedua.
Kita akan mempelajari pemetaan yang 
hanya diwajibkan mempertahankan struktur,
tanpa diwajibkan pula menjadi korespondensi. 

Pengalaman menunjukkan bahwa pemetaan-pemetaan ini 
sangat berguna.
Salah satunya, pada subbagian kedua di bawah ini kita akan melihat
bahwa isomorfisme menjelaskan bagaimana ruang-ruang itu sama,
sedangkan pemetaan-pemetaan ini dapat dipandang sebagai penjelasan tentang bagaimana ruang-ruang itu serupa.








\subsection{Definisi}

\begin{definition}  \label{def:Homo}
%<*df:Homo>
Suatu fungsi antarruang vektor \( \map{h}{V}{W} \) yang 
mempertahankan\index{mempertahankan struktur}\index{struktur!pelestarian} 
penjumlahan
\begin{center}
  jika \( \vec{v}_1,\vec{v}_2\in V \), maka
      \( h(\vec{v}_1+\vec{v}_2)=h(\vec{v}_1)+h(\vec{v}_2) \)
\end{center}
dan perkalian skalar
\begin{center}
      jika \( \vec{v}\in V \) dan \( r\in\Re \), maka
      \( h(r\cdot\vec{v})=r\cdot h(\vec{v}) \)
\end{center}
disebut \definend{homomorfisme}\index{homomorfisme}\index{pemetaan linear}%
\index{fungsi!mempertahankan struktur!\see{homomorfisme}}%
\index{ruang vektor!homomorfisme}\index{ruang vektor!pemetaan}
atau \definend{pemetaan linear}\index{pemetaan linear|seealso{homomorfisme}}.
%</df:Homo>
\end{definition}

\begin{example}    \label{ex:RThreeHomoRTwoFirst}
Pemetaan proyeksi\index{proyeksi}
\( \map{\pi}{\Re^3}{\Re^2} \)
\begin{equation*}
   \colvec{x \\ y \\ z}
    \mapsunder{\pi}
   \colvec{x \\ y}
\end{equation*}
adalah suatu homomorfisme.
Pemetaan ini mempertahankan penjumlahan
\begin{equation*}
  \pi(\colvec{x_1 \\ y_1 \\ z_1}\!+\!\colvec{x_2 \\ y_2 \\ z_2})
  =
  \pi(\colvec{x_1+x_2 \\ y_1+y_2 \\ z_1+z_2})
  =
  \colvec{x_1+x_2 \\ y_1+y_2}
  =
  \pi(\colvec{x_1 \\ y_1 \\ z_1})
  +
  \pi(\colvec{x_2 \\ y_2 \\ z_2})
\end{equation*}
dan perkalian skalar.
\begin{equation*}
  \pi(r\cdot\colvec{x_1 \\ y_1 \\ z_1})
  =
  \pi(\colvec{rx_1 \\ ry_1 \\ rz_1})
  =
  \colvec{rx_1 \\ ry_1}
  =
  r\cdot\pi(\colvec{x_1 \\ y_1 \\ z_1})
\end{equation*}
Pemetaan ini bukan isomorfisme karena tidak satu-ke-satu. 
Sebagai contoh, baik $\zero$ maupun $\vec{e}_3$ dalam $\Re^3$ dipetakan ke
vektor nol dalam $\Re^2$.
\end{example}

\begin{example} \label{exam:TwoMapsHomoNotIso}
Domain dan kodomain 
tidak harus berupa ruang vektor kolom.
Kedua pemetaan berikut adalah homomorfisme;
pembuktiannya langsung.
\begin{enumerate}
  \item \( \map{f_1}{\polyspace_2}{\polyspace_3} \) yang diberikan oleh
    \begin{equation*}
      a_0+a_1x+a_2x^2 \;\mapsto\; a_0x+(a_1/2)x^2+(a_2/3)x^3 
    \end{equation*}
  \item \( \map{f_2}{M_{\nbyn{2}}}{\Re} \) yang diberikan oleh
    \begin{equation*}
      \begin{mat}
        a  &b  \\
        c  &d
      \end{mat}
        \mapsto
      a+d
    \end{equation*}
\end{enumerate}
\end{example}

\begin{example}
Di antara sembarang dua ruang terdapat suatu \definend{homomorfisme nol},%
\index{homomorfisme nol}\index{homomorfisme!nol}\index{fungsi!nol}
yang memetakan setiap vektor dalam domain ke vektor nol dalam kodomain.
\end{example}

Kita akan menggunakan kedua istilah `homomorfisme' dan `pemetaan linear' secara bergantian.

\begin{example}
Dua contoh berikut memberi petunjuk mengapa kita menyebutnya `pemetaan linear'.
\begin{enumerate}
  \item Pemetaan \( \map{g}{\Re^3}{\Re} \) yang diberikan oleh
    \begin{equation*}
      \colvec{x \\ y \\ z}
        \mapsunder{g}
      3x+2y-4.5z
    \end{equation*}
    bersifat linear, yaitu merupakan homomorfisme.
    Pemeriksaannya mudah.
    Sebaliknya, pemetaan \( \map{\hat{g}}{\Re^3}{\Re} \) yang diberikan oleh
    \begin{equation*}
      \colvec{x \\ y \\ z}
        \mapsunder{\hat{g}}
      3x+2y-4.5z+1
    \end{equation*}
    tidak linear.
    Untuk menunjukkannya kita hanya perlu memberikan satu
    kombinasi linear yang tidak dipertahankan oleh pemetaan tersebut.
    Berikut salah satunya.
    \begin{equation*}
      \hat{g}(\colvec[r]{0 \\ 0 \\ 0}+\colvec[r]{1 \\ 0 \\ 0})=4
      \qquad
      \hat{g}(\colvec[r]{0 \\ 0 \\ 0})
      +\hat{g}(\colvec[r]{1 \\ 0 \\ 0})=5
    \end{equation*}
  \item Pemetaan pertama dari kedua pemetaan 
    \( \map{t_1,t_2}{\Re^3}{\Re^2} \)
    berikut bersifat linear, sedangkan yang kedua tidak.
    \begin{equation*}
      \colvec{x \\ y \\ z}
        \mapsunder{t_1}
      \colvec{5x-2y \\ x+y}
      \qquad
      \colvec{x \\ y \\ z}
        \mapsunder{t_2}
      \colvec{5x-2y \\ xy}
    \end{equation*}
    Mudah untuk menemukan kombinasi linear yang tidak 
    dipertahankan oleh pemetaan kedua. 
\end{enumerate}
% The homomorphisms have 
% coordinate functions that are linear combinations of the arguments.
% See also \nearbyexercise{exer:GrpahNotALine}.
\end{example}

Jadi, salah satu cara memandang `homomorfisme' 
adalah bahwa kita sedang memperumum `isomorfisme' (dengan membuang syarat bahwa
pemetaan tersebut merupakan korespondensi),
berdasarkan pengamatan bahwa banyak sifat
isomorfisme hanya bergantung pada sifat pemetaan yang mempertahankan struktur.
Dua hasil berikut merupakan contoh dari gagasan ini.
Pada bagian sebelumnya kita melihat bahwa bukti untuk masing-masing hasil itu hanya menggunakan pelestarian
penjumlahan dan pelestarian perkalian skalar,
sehingga berlaku pula untuk homomorfisme.

\begin{lemma}       \label{le:HomoSendsZeroToZero}
%<*lm:HomoSendsZeroToZero>
Pemetaan linear memetakan vektor nol ke vektor nol.
%</lm:HomoSendsZeroToZero>
\end{lemma}

\begin{lemma}  \label{le:HomoPreserveLinCombo}
%<*lm:HomoPreserveLinCombo>
Pernyataan-pernyataan berikut ekuivalen untuk sembarang pemetaan 
\( \map{f}{V}{W} \) 
antarruang vektor.
\begin{tfae}
  \item 
      $f$ adalah homomorfisme 
  \item 
      $f(c_1\cdot\vec{v}_1+c_2\cdot\vec{v}_2)
      =c_1\cdot f(\vec{v}_1)+c_2\cdot f(\vec{v}_2)$
      untuk sembarang \( c_1,c_2\in\Re \) dan \( \vec{v}_1,\vec{v}_2\in V \)
  \item
    $f(c_1\cdot\vec{v}_1+\dots+c_n\cdot\vec{v}_n)
    =c_1\cdot f(\vec{v}_1)+\dots+c_n\cdot f(\vec{v}_n)$ 
    untuk sembarang \( c_1,\dots,c_n\in\Re \) dan
    \( \vec{v}_1,\ldots,\vec{v}_n\in V \)
\end{tfae}
%</lm:HomoPreserveLinCombo>
\end{lemma}

\begin{example}
Fungsi \( \map{f}{\Re^2}{\Re^4} \) yang diberikan oleh
\begin{equation*}
  \colvec{x \\ y}
    \mapsunder{f}
  \colvec{x/2 \\ 0 \\ x+y \\ 3y}
\end{equation*}
bersifat linear karena memenuhi butir~(2).
\begin{equation*}
  \colvec{r_1(x_1/2)+r_2(x_2/2) \\ 0 \\ 
                 r_1(x_1+y_1)+r_2(x_2+y_2) \\ r_1(3y_1)+r_2(3y_2)}
   =
  r_1\colvec{x_1/2 \\ 0 \\ x_1+y_1 \\ 3y_1}
   +
  r_2\colvec{x_2/2 \\ 0 \\ x_2+y_2 \\ 3y_2}
\end{equation*}
\end{example}

Namun,
sebagian hal yang berlaku untuk isomorfisme tidak berlaku untuk
homomorfisme.
Salah satu contohnya terdapat dalam bukti Lema~I.\ref{lem:IsoImpliesSameDim}, 
yang menunjukkan bahwa isomorfisme antarruang menghasilkan 
korespondensi di antara basis-basisnya.
Homomorfisme tidak menghasilkan korespondensi semacam itu;
\nearbyexample{ex:RThreeHomoRTwoFirst} menunjukkan hal ini, dan contoh lainnya adalah
pemetaan nol di antara dua ruang nontrivial. 
Sebagai gantinya, untuk homomorfisme kita memiliki hasil yang lebih lemah tetapi tetap sangat berguna.

\begin{theorem} \label{th:HomoDetActOnBasis}
%<*th:HomoDetActOnBasis>
Suatu homomorfisme ditentukan oleh aksinya pada suatu basis:~jika
$V$ adalah ruang vektor dengan basis
\( \sequence{\vec{\beta}_1,\dots,\vec{\beta}_n} \),
Jika $W$ adalah ruang vektor, dan jika 
\( \vec{w}_1,\dots,\vec{w}_n\in W \) 
(unsur-unsur kodomain ini tidak harus berbeda), maka
terdapat suatu homomorfisme dari \( V \) ke \( W \) yang memetakan setiap
\( \vec{\beta}_i \) ke \( \vec{w}_i \), dan homomorfisme tersebut tunggal.
%</th:HomoDetActOnBasis>
\end{theorem}

\begin{proof}
%<*pf:HomoDetActOnBasis0>
Untuk sembarang masukan $\vec{v}\in V$, misalkan bentuknya terhadap
basis tersebut adalah
\( \vec{v}=c_1\vec{\beta}_1+\dots+c_n\vec{\beta}_n \). 
Definisikan
keluaran yang terkait dengan menggunakan koordinat yang sama,
$h(\vec{v})=c_1\vec{w}_1+\dots+c_n\vec{w}_n$.
Fungsi ini terdefinisi dengan baik karena, terhadap basis tersebut,
representasi setiap vektor domain \( \vec{v} \) bersifat tunggal.
%</pf:HomoDetActOnBasis0>

%<*pf:HomoDetActOnBasis1>
Pemetaan ini merupakan homomorfisme 
karena mempertahankan kombinasi linear:
dengan \( \vec{v_1}=c_1\vec{\beta}_1+\cdots+c_n\vec{\beta}_n \) dan
\( \vec{v_2}=d_1\vec{\beta}_1+\cdots+d_n\vec{\beta}_n \), berikut
perhitungannya.
\begin{align*}
  h(r_1\vec{v}_1+r_2\vec{v}_2)
  &=h(\,(r_1c_1+r_2d_1)\vec{\beta}_1+\dots+(r_1c_n+r_2d_n)\vec{\beta}_n\,)  \\
  &=(r_1c_1+r_2d_1)\vec{w}_1+\dots+(r_1c_n+r_2d_n)\vec{w}_n   \\
  &=r_1h(\vec{v}_1)+r_2h(\vec{v}_2)
\end{align*}
%</pf:HomoDetActOnBasis1>

%<*pf:HomoDetActOnBasis2>
Pemetaan ini tunggal karena jika \( \map{\hat{h}}{V}{W} \) 
adalah homomorfisme lain yang memenuhi \( \hat{h}(\vec{\beta}_i)=\vec{w}_i \) 
untuk setiap \( i \),
maka \( h \) dan \( \hat{h} \) memberikan 
hasil yang sama pada semua vektor dalam domain. 
\begin{multline*}
  \hat{h}(\vec{v})
  =\hat{h}(c_1\vec{\beta}_1+\dots+c_n\vec{\beta}_n)  
  =c_1 \hat{h}(\vec{\beta}_1)+\dots+c_n \hat{h}(\vec{\beta}_n)  \\  
  =c_1\vec{w}_1+\dots+c_n\vec{w}_n 
  =h(\vec{v})
\end{multline*}
Keduanya memiliki aksi yang sama, sehingga merupakan fungsi yang sama.
%</pf:HomoDetActOnBasis2>
\end{proof}

\begin{definition} \label{df:ExtendedLinearly}
%<*df:ExtendedLinearly>
Misalkan $V$ dan~$W$ adalah ruang vektor dan 
misalkan  
$B=\sequence{\vec{\beta}_1,\ldots,\vec{\beta}_n}$ 
adalah basis untuk~$V$.  
Suatu fungsi yang didefinisikan pada basis tersebut $\map{f}{B}{W}$
\definend{diperluas secara linear}\index{diperluas, secara linear}\index{fungsi!diperluas secara linear}\index{perluasan linear suatu fungsi}
menjadi fungsi $\map{\hat{f}}{V}{W}$ jika,
untuk setiap $\vec{v}\in V$ sedemikian sehingga
$\vec{v}=c_1\vec{\beta}_1+\cdots+c_n\vec{\beta}_n$,
aksi pemetaannya adalah
$\hat{f}(\vec{v})=c_1\cdot f(\vec{\beta}_1)
  +\cdots+c_n\cdot f(\vec{\beta}_n)$.
%</df:ExtendedLinearly>
\end{definition}


\begin{example}
Jika kita menentukan suatu pemetaan \( \map{h}{\Re^2}{\Re^2} \)
yang beraksi pada basis standar $\stdbasis_2$ dengan cara berikut, 
\begin{equation*}
  h(\colvec[r]{1 \\ 0})=\colvec[r]{-1 \\ 1}
  \qquad
  h(\colvec[r]{0 \\ 1})=\colvec[r]{-4 \\ 4}
\end{equation*}
maka kita juga telah menentukan aksi $h$ pada setiap anggota domain lainnya.
Sebagai contoh,
nilai $h$ pada argumen berikut
\begin{equation*}
  h(\colvec[r]{3 \\ -2})=h(3\cdot \colvec[r]{1 \\ 0}-2\cdot \colvec[r]{0 \\ 1})
                      =3\cdot h(\colvec[r]{1 \\ 0})-2\cdot h(\colvec[r]{0 \\ 1})
                      =\colvec[r]{5 \\ -5}
\end{equation*}
merupakan konsekuensi langsung dari nilai $h$ pada vektor-vektor basis.
\end{example}

Nanti dalam bab ini kita akan mengembangkan suatu skema yang praktis untuk perhitungan 
seperti ini dengan menggunakan matriks.

% Just as the isomorphisms of a space with itself are useful and interesting, 
% so too are the homomorphisms of a space with itself.

\begin{definition} \label{df:LinearTransformation}
%<*df:LinearTransformation>
Pemetaan linear dari suatu ruang ke ruang itu sendiri \( \map{t}{V}{V} \) disebut
\definend{transformasi linear}\index{transformasi linear}\index{transformasi linear|seealso{transformasi}}.
%</df:LinearTransformation>
\end{definition}

\begin{remark}
Dalam buku ini kita menggunakan `transformasi linear' hanya ketika 
kodomain sama dengan domain.
Perhatikan bahwa sebagian sumber justru menggunakannya sebagai sinonim `pemetaan linear'.
Sinonim lainnya adalah `operator linear'.
\end{remark}

\begin{example}
Pemetaan pada $\Re^2$ yang memproyeksikan semua vektor ke sumbu-$x$
merupakan transformasi linear.
\begin{equation*}
  \colvec{x \\ y}\mapsto\colvec{x \\ 0}
\end{equation*}
\end{example}

\begin{example}
Pemetaan turunan \( \map{d/dx}{\polyspace_n}{\polyspace_n} \)
\begin{equation*}
  a_0+a_1x+\cdots+a_nx^n
    \mapsunder{d/dx}
  a_1+2a_2x+3a_3x^2+\cdots+na_nx^{n-1}
\end{equation*}
merupakan transformasi linear, sebagaimana ditunjukkan oleh hasil dari kalkulus berikut:
\( d(c_1f+c_2g)/dx=c_1\,(df/dx)+c_2\,(dg/dx) \).
\end{example}

\begin{example} \label{ex:MatTransMapLinear}
Operasi transpos matriks
\begin{equation*}
  \begin{mat}
    a  &b  \\
    c  &d
  \end{mat}
  \;\mapsto\;
  \begin{mat}
    a  &c  \\
    b  &d
  \end{mat}
\end{equation*}
merupakan transformasi linear pada \( \matspace_{\nbyn{2}} \).
(Transpos bersifat satu-ke-satu dan pada, sehingga sebenarnya 
merupakan automorfisme.)
\end{example}

Kita mengakhiri subbagian tentang pemetaan ini dengan mengingat kembali bahwa
kita dapat mengombinasikan pemetaan secara linear.
Sebagai contoh, untuk pemetaan-pemetaan berikut dari \( \Re^2 \) ke dirinya sendiri,
\begin{equation*}
  \colvec{x \\ y}
   \mapsunder{f}
  \colvec{2x \\ 3x-2y}
  \quad\text{dan}\quad
  \colvec{x \\ y}
   \mapsunder{g}
  \colvec{0 \\ 5x}
\end{equation*}
kombinasi linear \( 5f-2g \) juga merupakan transformasi pada~$\Re^2$.
\begin{equation*}
  \colvec{x \\ y}
   \mapsunder{5f-2g}
  \colvec{10x \\ 5x-10y}
\end{equation*}

\begin{lemma} \label{le:SpLinFcns}
%<*lm:SpLinFcns>
Untuk ruang vektor \( V \) dan \( W \),
himpunan fungsi linear dari \( V \) ke
\( W \) itu sendiri merupakan ruang vektor, yakni subruang dari ruang semua fungsi
dari \( V \) ke \( W \).
%</lm:SpLinFcns>
\end{lemma}

%<*SpLinFcns>
\noindent Kita menyatakan ruang pemetaan linear dari $V$ ke~$W$ dengan
\( \linmaps{V}{W} \).\index{pemetaan linear, ruang vektor dari}
%</SpLinFcns>

\begin{proof}
%<*pf:SpLinFcns>
Himpunan ini tidak kosong karena memuat homomorfisme nol.
Jadi, untuk menunjukkan bahwa himpunan ini merupakan subruang, kita hanya perlu memeriksa bahwa himpunan ini 
tertutup terhadap operasi-operasinya.
Misalkan \( \map{f,g}{V}{W} \) linear. 
Maka operasi penjumlahan fungsi dipertahankan,
\begin{align*}
   (f+g)(c_1\vec{v}_1+c_2\vec{v}_2)
   &=f(c_1\vec{v}_1+c_2\vec{v}_2) + 
   g(c_1\vec{v}_1+c_2\vec{v}_2)       \\
   &=c_1f(\vec{v}_1)+c_2f(\vec{v}_2)
   +c_1g(\vec{v}_1)+c_2g(\vec{v}_2)   \\
   &=c_1\bigl(f+g\bigr)(\vec{v}_1)+c_2\bigl(f+g\bigr)(\vec{v}_2)
\end{align*}
demikian pula operasi perkalian skalar pada suatu fungsi.
\begin{align*}
   (r\cdot f)(c_1\vec{v}_1+c_2\vec{v}_2)
   &=r(c_1f(\vec{v}_1)+c_2f(\vec{v}_2))  \\
   &=c_1(r\cdot f)(\vec{v}_1)+c_2(r\cdot f)(\vec{v}_2)
\end{align*}
Dengan demikian, \( \linmaps{V}{W} \) merupakan subruang.
%</pf:SpLinFcns>
\end{proof}

Kita memulai bagian ini dengan 
mendefinisikan `homomorfisme' sebagai perumuman `isomorfisme',
dengan memisahkan sifat pelestarian struktur.
Sebagian hal mengenai isomorfisme berlaku tanpa perubahan, sedangkan
hal lainnya kita sesuaikan.

Namun, perhatikan bahwa gagasan 
`homomorfisme' sama sekali bukan gagasan sekunder dibandingkan
dengan `isomorfisme'.
Pada sisa bab ini kita terutama akan bekerja dengan homomorfisme.
Hal ini
sebagian karena setiap pernyataan tentang homomorfisme secara otomatis berlaku 
untuk isomorfisme, tetapi terutama karena, 
meskipun konsep isomorfisme terasa lebih alami, 
pengalaman kita akan menunjukkan bahwa konsep homomorfisme 
lebih produktif dan lebih sentral bagi kemajuan kita.

\begin{exercises}
  \recommended \item
    Tentukan apakah setiap \( \map{h}{\Re^3}{\Re^2} \) berikut linear.
    \begin{exparts*}
      \partsitem \( h(\colvec{x \\ y \\ z})=\colvec{x \\ x+y+z}  \)
      \partsitem \( h(\colvec{x \\ y \\ z})=\colvec[r]{0 \\ 0}  \)
      \partsitem \( h(\colvec{x \\ y \\ z})=\colvec[r]{1 \\ 1}  \)
      \partsitem \( h(\colvec{x \\ y \\ z})=\colvec{2x+y \\ 3y-4z}  \)
    \end{exparts*}
    \begin{answer}
      \begin{exparts}
        \partsitem Ya.
          Pembuktiannya langsung.
          \begin{align*}
            h( c_1\cdot\colvec{x_1 \\ y_1 \\ z_1}
               +c_2\cdot\colvec{x_2 \\ y_2 \\ z_2} )
            &=h( \colvec{c_1x_1+c_2x_2 \\ c_1y_1+c_2y_2 \\ c_1z_1+c_2z_2} ) \\
            &=\colvec{c_1x_1+c_2x_2 \\ 
                     c_1x_1+c_2x_2+c_1y_1+c_2y_2+c_1z_1+c_2z_2}          \\
            &=c_1\cdot\colvec{x_1 \\ x_1+y_1+z_1}
               +c_2\cdot\colvec{x_2 \\ x_2+y_2+z_2}                       \\  
            &=c_1\cdot h(\colvec{x_1 \\ y_1 \\ z_1}) 
               +c_2\cdot h(\colvec{x_2 \\ y_2 \\ z_2})
          \end{align*}
        \partsitem Ya.
          Pembuktiannya mudah.
          \begin{align*}
            h(c_1\cdot\colvec{x_1 \\ y_1 \\ z_1}
              +c_2\cdot\colvec{x_2 \\ y_2 \\ z_2})
            &=h( \colvec{c_1x_1+c_2x_2 \\ c_1y_1+c_2y_2 \\ c_1z_1+c_2z_2} ) \\
            &=\colvec[r]{0 \\ 0}                                      \\
            &=c_1\cdot h(\colvec{x_1 \\ y_1 \\ z_1}) 
               +c_2\cdot h(\colvec{x_2 \\ y_2 \\ z_2})            
          \end{align*}
        \partsitem Tidak.
          Berikut contoh penjumlahan yang tidak dipertahankan.
          \begin{equation*}
            h(\colvec[r]{0 \\ 0 \\ 0}+\colvec[r]{0 \\ 0 \\ 0})
            =\colvec[r]{1 \\ 1}
            \neq h(\colvec[r]{0 \\ 0 \\ 0})+h(\colvec[r]{0 \\ 0 \\ 0})
          \end{equation*}
        \partsitem Ya.
           Pembuktiannya langsung.
          \begin{align*}
            h( c_1\cdot\colvec{x_1 \\ y_1 \\ z_1}
               +c_2\cdot\colvec{x_2 \\ y_2 \\ z_2} )
            &=h( \colvec{c_1x_1+c_2x_2 \\ c_1y_1+c_2y_2 \\ c_1z_1+c_2z_2} ) \\
            &=\colvec{2(c_1x_1+c_2x_2)+(c_1y_1+c_2y_2) \\ 
                      3(c_1y_1+c_2y_2)-4(c_1z_1+c_2z_2)}          \\
            &=c_1\cdot\colvec{2x_1+y_1 \\ 3y_1-4z_1}
               +c_2\cdot\colvec{2x_2+y_2 \\ 3y_2-4z_2}            \\  
            &=c_1\cdot h(\colvec{x_1 \\ y_1 \\ z_1}) 
               +c_2\cdot h(\colvec{x_2 \\ y_2 \\ z_2})
          \end{align*}
      \end{exparts}  
    \end{answer}
  \recommended \item
    Tentukan apakah setiap pemetaan \( \map{h}{\matspace_{\nbyn{2}}}{\Re} \) berikut linear.
    \begin{exparts}
      \partsitem \( h(\begin{mat} a &b  \\ c &d \end{mat})=a+d  \)
      \partsitem \( h(\begin{mat} a &b  \\ c &d \end{mat})=ad-bc  \)
      \partsitem \( h(\begin{mat} a &b \\ c &d \end{mat})=2a+3b+c-d  \)
      \partsitem \( h(\begin{mat} a &b  \\ c &d \end{mat})=a^2+b^2  \)
    \end{exparts}
    \begin{answer} 
      Untuk masing-masing pemetaan, kita harus memeriksa bahwa pemetaan itu mempertahankan kombinasi linear
      atau memberikan contoh kombinasi linear yang tidak 
      dipertahankannya.
      \begin{exparts*}
        \partsitem Ya.
           Pemeriksaan bahwa pemetaan ini mempertahankan kombinasi bersifat rutin.
           \begin{align*}
             h(r_1\cdot\begin{mat}
                 a_1  &b_1  \\
                 c_1  &d_1
               \end{mat}
              +r_2\cdot\begin{mat}
                 a_2  &b_2  \\
                 c_2  &d_2
               \end{mat})
              &=h(\begin{mat}
                 r_1a_1+r_2a_2  &r_1b_1+r_2b_2  \\
                 r_1c_1+r_2c_2  &r_1d_1+r_2d_2
               \end{mat})                    \\
              &=(r_1a_1+r_2a_2)+(r_1d_1+r_2d_2)    \\
              &=r_1(a_1+d_1)+r_2(a_2+d_2)          \\
              &=r_1\cdot h(\begin{mat}
                             a_1  &b_1  \\
                             c_1  &d_1
                            \end{mat})
                +r_2\cdot h(\begin{mat}
                             a_2  &b_2  \\
                             c_2  &d_2
                            \end{mat})
           \end{align*}
        \partsitem Tidak.
          Sebagai contoh, perkalian dengan skalar $2$ tidak dipertahankan.
          \begin{equation*}
            h(2\cdot\begin{mat}[r]
                1  &0  \\
                0  &1  
              \end{mat})
            =h(\begin{mat}[r]
                2  &0  \\
                0  &2  
              \end{mat})
            =4
            \quad\text{sedangkan}\quad
            2\cdot h(\begin{mat}[r]
                1  &0  \\
                0  &1  
              \end{mat})
            =2\cdot 1=2
          \end{equation*}
        \partsitem Ya.
           Berikut pemeriksaan bahwa pemetaan ini mempertahankan kombinasi dua anggota
           domain.
           \begin{multline*}
             h(r_1\cdot\begin{mat} a_1 &b_1 \\ c_1 &d_1 \end{mat}
               +r_2\cdot\begin{mat} a_2 &b_2 \\ c_2 &d_2 \end{mat})     \\
             \begin{aligned}
             &=h(\begin{mat} 
                   r_1a_1+r_2a_2 &r_1b_1+r_2b_2   \\ 
                   r_1c_1+r_2c_2 &r_1d_1+r_2d_2 
                 \end{mat})                                 \\
             &=2(r_1a_1+r_2a_2)+3(r_1b_1+r_2b_2)
                    +(r_1c_1+r_2c_2)-(r_1d_1+r_2d_2)              \\
             &=r_1(2a_1+3b_1+c_1-d_1)
               +r_2(2a_2+3b_2+c_2-d_2)              \\
             &=r_1\cdot h(\begin{mat} a_1 &b_1 \\ c_1 &d_1 \end{mat}
               +r_2\cdot h(\begin{mat} a_2 &b_2 \\ c_2 &d_2 \end{mat})
             \end{aligned}
           \end{multline*}
        \partsitem Tidak.
          Sebagai contoh kombinasi yang tidak dipertahankan,
          menghitungnya dengan satu cara memberikan hasil berikut,
          \begin{equation*}
            h(\begin{mat}[r]
              1  &0  \\
              0  &0
            \end{mat}
            +\begin{mat}[r]
              1  &0  \\
              0  &0
            \end{mat})
           =h(\begin{mat}[r]
              2  &0  \\
              0  &0
            \end{mat})
           =4
           \end{equation*}
           sedangkan cara lainnya memberikan hasil berikut.
           \begin{equation*}
            h(\begin{mat}[r]
              1  &0  \\
              0  &0
            \end{mat})
            +h(\begin{mat}[r]
              1  &0  \\
              0  &0
            \end{mat})
            =1+1
            =2
          \end{equation*}
      \end{exparts*}   
    \end{answer}
  \recommended \item
    Tunjukkan bahwa pemetaan-pemetaan berikut adalah homomorfisme.
    Apakah keduanya saling invers?
    \begin{exparts}
      \partsitem \( \map{d/dx}{\polyspace_3}{\polyspace_2} \)
        yang diberikan dengan memetakan \( a_0+a_1x+a_2x^2+a_3x^3 \) ke
        \( a_1+2a_2x+3a_3x^2 \)
      \partsitem \( \map{\int}{\polyspace_2}{\polyspace_3} \) yang diberikan dengan memetakan
        \( b_0+b_1x+b_2x^2 \) ke \( b_0x+(b_1/2)x^2+(b_2/3)x^3 \)
    \end{exparts}
    \begin{answer}
      Pemeriksaan bahwa masing-masing merupakan homomorfisme bersifat rutin.
      Berikut pemeriksaan untuk pemetaan diferensiasi.
      \begin{multline*}
        \frac{d}{dx}(r\cdot (a_0+a_1x+a_2x^2+a_3x^3)
                    +s\cdot (b_0+b_1x+b_2x^2+b_3x^3))                 \\
        \begin{aligned}
           &=\frac{d}{dx}((ra_0+sb_0)+(ra_1+sb_1)x+(ra_2+sb_2)x^2
                                                    +(ra_3+sb_3)x^3)  \\
           &=(ra_1+sb_1)+2(ra_2+sb_2)x+3(ra_3+sb_3)x^2         \\
           &=r\cdot (a_1+2a_2x+3a_3x^2)+s\cdot (b_1+2b_2x+3b_3x^2)         \\
           &=r\cdot \frac{d}{dx}(a_0+a_1x+a_2x^2+a_3x^3)
                    +s\cdot \frac{d}{dx} (b_0+b_1x+b_2x^2+b_3x^3)
        \end{aligned}
      \end{multline*}
      (Bukti alternatifnya cukup dengan memperhatikan bahwa ini merupakan 
      sifat diferensiasi yang sudah dikenal dari kalkulus.)

      Kedua pemetaan ini bukan saling invers karena komposisi berikut 
      tidak beraksi sebagai pemetaan identitas pada 
      unsur domain ini.
      \begin{equation*}
        1\in\polyspace_3\;\mapsunder{d/dx}\;
        0\in\polyspace_2\;\mapsunder{\int}\;
        0\in\polyspace_3
      \end{equation*}   
     \end{answer}
  \item  
    Apakah proyeksi (tegak lurus) dari \( \Re^3 \) ke bidang-\( xz \)
    merupakan homomorfisme?
    Bagaimana dengan proyeksi ke bidang-\( yz \)?
    Ke sumbu-\( x \)?
    Sumbu-\( y \)?
    Sumbu-\( z \)?
    Proyeksi ke titik asal?
    \begin{answer}
      Setiap proyeksi tersebut merupakan homomorfisme.
      Proyeksi ke bidang-$xz$ dan ke bidang-$yz$ adalah pemetaan berikut.
      \begin{equation*}
         \colvec{x \\ y \\ z}\mapsto\colvec{x \\ 0 \\ z}
           \qquad
         \colvec{x \\ y \\ z}\mapsto\colvec{0 \\ y \\ z}
      \end{equation*}
      Proyeksi ke sumbu-$x$, ke sumbu-$y$, dan ke sumbu-$z$ adalah
      pemetaan berikut. 
      \begin{equation*}
         \colvec{x \\ y \\ z}\mapsto\colvec{x \\ 0 \\ 0}
           \qquad
         \colvec{x \\ y \\ z}\mapsto\colvec{0 \\ y \\ 0}
           \qquad
         \colvec{x \\ y \\ z}\mapsto\colvec{0 \\ 0 \\ z}
      \end{equation*}
      Sementara itu, proyeksi ke titik asal adalah pemetaan berikut.
      \begin{equation*}
         \colvec{x \\ y \\ z}\mapsto\colvec[r]{0 \\ 0 \\ 0}
      \end{equation*}
      Pembuktian bahwa masing-masing merupakan homomorfisme bersifat langsung.
      (Yang terakhir, tentu saja, adalah transformasi nol pada $\Re^3$.)  
     \end{answer}
  \item Buktikan bahwa setiap pemetaan berikut merupakan homomorfisme.
    \begin{exparts}
      \partsitem $\map{h}{\polyspace_2}{\Re^2}$ yang diberikan oleh
        \begin{equation*}
          ax^2+bx+c\mapsto \colvec{a+b \\ a+c}
        \end{equation*}
      \partsitem $\map{f}{\Re^2}{\Re^3}$ yang diberikan oleh 
        \begin{equation*}
          \colvec{x \\ y}\mapsto\colvec{0 \\ x-y \\ 3y} 
        \end{equation*}
    \end{exparts}
    \begin{answer}
      \begin{exparts}
        \partsitem
          Perhitungan ini membuktikan bahwa pemetaan tersebut mempertahankan kombinasi linear.
          Menurut \nearbylemma{le:HomoPreserveLinCombo}, hal ini cukup untuk menunjukkan bahwa 
          pemetaan tersebut merupakan homomorfisme.
          \begin{multline*}
            h(\,d_1(a_1x^2+b_1x+c_1)+d_2(a_2x^2+b_2x+c_2)\,)      \\
             \begin{aligned}
             &=h((d_1a_1+d_2a_2)x^2+(d_1b_1+d_2b_2)x+(d_1c_1+d_2c_2))  \\
             &=\colvec{(d_1a_1+d_2a_2)+(d_1b_1+d_2b_2) \\ 
             (d_1a_1+d_2a_2)+(d_1c_1+d_2c_2)} \\
             &=\colvec{d_1a_1+d_1b_1 \\  d_1a_1+d_1c_1}+
                 \colvec{d_2a_2+d_2b_2 \\  d_2a_2+d_2c_2}  \\
             &=d_1\colvec{a_1+b_1 \\  a_1+c_1}+
                 d_2\colvec{a_2+b_2 \\  a_2+c_2}  \\
            &=d_1\cdot h(a_1x^2+b_1x+c_1)+d_2\cdot h(a_2x^2+b_2x+c_2)
            \end{aligned}
          \end{multline*}
        \partsitem
          Pemetaan ini mempertahankan kombinasi linear.
          \begin{align*}
            f(\,a_1\colvec{x_1 \\ y_1}+a_2\colvec{x_2 \\ y_2}\,)
              &=f(\,\colvec{a_1x_1+a_2x_2 \\ a_1y_1+a_2y_2}\,)  \\
              &=\colvec{0 \\ (a_1x_1+a_2x_2)-(a_1y_1+a_2y_2) \\ 
              3(a_1y_1+a_2y_2) }  \\
              &=a_1\colvec{0 \\ x_1-y_1 \\ 3y_1}
                 +a_2\colvec{0 \\ x_2-y_2 \\ 3y_2}  \\
              &=a_1f(\colvec{x_1 \\ y_1})+a_2f(\colvec{x_2 \\ y_2})
          \end{align*}
      \end{exparts}
    \end{answer}


  \item 
    Tunjukkan bahwa, meskipun pemetaan-pemetaan dari \nearbyexample{exam:TwoMapsHomoNotIso}
    mempertahankan operasi linear, keduanya bukan isomorfisme.
    \begin{answer}
      Pemetaan pertama tidak pada; sebagai contoh, tidak ada polinomial yang
      dipetakan ke polinomial konstan $p(x)=1$.
      Pemetaan kedua tidak satu-ke-satu; kedua anggota domain berikut
      \begin{equation*}
        \begin{mat}[r]
          1  &0  \\
          0  &0  
        \end{mat}
        \quad\text{dan}\quad
        \begin{mat}[r]
          0  &0  \\
          0  &1  
        \end{mat}
      \end{equation*}
      dipetakan ke anggota kodomain yang sama, yaitu $1\in\Re$.
     \end{answer}
  \item 
    Apakah pemetaan identitas merupakan transformasi linear?
    \begin{answer}
      Ya;~dalam sembarang ruang, \( \text{id}(c\cdot \vec{v}+d\cdot \vec{w})
      = c\cdot \vec{v}+d\cdot \vec{w}
      = c\cdot\text{id}(\vec{v})+d\cdot\text{id}(\vec{w}) \).
    \end{answer}
  \recommended \item \label{exer:GrpahNotALine}
    Menyatakan bahwa suatu fungsi `linear' berbeda dengan 
    menyatakan bahwa grafiknya berupa garis.
    \begin{exparts}
      \partsitem Fungsi \( \map{f_1}{\Re}{\Re} \) yang diberikan oleh 
        \( f_1(x)=2x-1 \) memiliki grafik berupa garis.
        Tunjukkan bahwa fungsi ini bukan fungsi linear.
      \partsitem Fungsi \( \map{f_2}{\Re^2}{\Re} \) yang diberikan oleh
        \begin{equation*}
          \colvec{x \\ y} \mapsto x+2y
        \end{equation*}
        tidak memiliki grafik berupa garis.
        Tunjukkan bahwa fungsi ini merupakan fungsi linear.
    \end{exparts}
    \begin{answer}
      \begin{exparts}
         \partsitem Pemetaan ini tidak mempertahankan struktur karena
           \( f(1+1)=3 \), sedangkan \( f(1)+f(1)=2 \).
         \partsitem Pemeriksaannya rutin.
           \begin{align*}
             f(r_1\cdot\colvec{x_1 \\ y_1}+r_2\cdot\colvec{x_2 \\ y_2})
             &=f(\colvec{r_1x_1+r_2x_2 \\ r_1y_1+r_2y_2})               \\
             &=(r_1x_1+r_2x_2)+2(r_1y_1+r_2y_2)                          \\
             &=r_1\cdot (x_1+2y_1)+r_2\cdot (x_2+2y_2)                   \\
             &=r_1\cdot f(\colvec{x_1 \\ y_1})+r_2\cdot f(\colvec{x_2 \\ y_2})
           \end{align*}
      \end{exparts}   
     \end{answer}
  \recommended \item 
    Salah satu bagian definisi fungsi linear adalah bahwa fungsi tersebut mempertahankan
    penjumlahan. 
    Apakah fungsi linear mempertahankan pengurangan?
    \begin{answer}
      Ya.
      Jika \( \map{h}{V}{W} \) linear, maka
      \( h(\vec{u}-\vec{v})
         =h(\vec{u}+(-1)\cdot\vec{v})
         =h(\vec{u})+(-1)\cdot h(\vec{v})
         =h(\vec{u})-h(\vec{v}) \).
    \end{answer}
  \item 
    Anggaplah \( h \) merupakan transformasi linear pada \( V \) dan
    \( \sequence{\vec{\beta}_1,\ldots,\vec{\beta}_n} \) merupakan basis \( V \).
    Buktikan setiap pernyataan berikut.
    \begin{exparts}
      \partsitem Jika \( h(\vec{\beta}_i)=\zero \) untuk setiap vektor basis,
        maka \( h \) adalah pemetaan nol.
      \partsitem Jika \( h(\vec{\beta}_i)=\vec{\beta}_i \) untuk setiap vektor
        basis, maka \( h \) adalah pemetaan identitas.
      \partsitem Jika terdapat skalar \( r \) sedemikian sehingga
        \( h(\vec{\beta}_i)=r\cdot\vec{\beta}_i \) untuk setiap
        vektor basis, maka \( h(\vec{v})=r\cdot\vec{v} \) untuk semua vektor
        dalam $V$.
    \end{exparts}
    \begin{answer}
      \begin{exparts}
       \partsitem Misalkan \( \vec{v}\in V \) direpresentasikan terhadap 
         basis tersebut sebagai \( \vec{v}=c_1\vec{\beta}_1+\dots+c_n\vec{\beta}_n \).
         Maka \( h(\vec{v})=h(c_1\vec{\beta}_1+\dots+c_n\vec{\beta}_n)
                  =c_1h(\vec{\beta}_1)+\dots+c_nh(\vec{\beta}_n)
                  =c_1\cdot\zero+\dots+c_n\cdot\zero
                  =\zero \).
       \partsitem Argumen ini serupa dengan argumen sebelumnya.
         Misalkan \( \vec{v}\in V \) direpresentasikan terhadap 
         basis tersebut sebagai \( \vec{v}=c_1\vec{\beta}_1+\dots+c_n\vec{\beta}_n \).
         Maka \( h(c_1\vec{\beta}_1+\dots+c_n\vec{\beta}_n)
            =c_1h(\vec{\beta}_1)+\dots+c_nh(\vec{\beta}_n)
            =c_1\vec{\beta}_1+\dots+c_n\vec{\beta}_n
            =\vec{v} \).
       \partsitem Seperti di atas, hanya saja 
         \( c_1h(\vec{\beta}_1)+\dots+c_nh(\vec{\beta}_n)
            =c_1r\vec{\beta}_1+\dots+c_nr\vec{\beta}_n
            =r(c_1\vec{\beta}_1+\dots+c_n\vec{\beta}_n)
            =r\vec{v} \).
      \end{exparts}  
     \end{answer}
  \item
    Tinjau ruang vektor \( \Re^+ \), dengan penjumlahan vektor
    dan perkalian skalar yang bukan diwarisi dari $\Re$,
    melainkan didefinisikan sebagai berikut:
    \( a+b \) adalah hasil kali \( a \) dan \( b \), sedangkan \( r\cdot a \) adalah
    pangkat ke-\( r \) dari \( a \).
    (Dalam latihan sebelumnya telah ditunjukkan bahwa ini merupakan ruang vektor.)
    Buktikan bahwa pemetaan logaritma alami \( \map{\ln}{\Re^+}{\Re} \) 
    merupakan homomorfisme di antara kedua ruang tersebut.
    Apakah pemetaan ini merupakan isomorfisme?
    \begin{answer}
      Fakta bahwa pemetaan ini merupakan homomorfisme mengikuti aturan yang sudah dikenal, yaitu
      logaritma suatu hasil kali adalah jumlah logaritmanya,
      $\ln(ab)=\ln(a)+\ln(b)$ 
      dan logaritma suatu pangkat adalah kelipatan logaritmanya,
      $\ln(a^r)=r\ln(a)$.
      Pemetaan ini merupakan isomorfisme karena memiliki invers, yaitu 
      pemetaan eksponensial, sehingga merupakan korespondensi,
      dan dengan demikian merupakan isomorfisme.  
     \end{answer}
  \item 
     Tinjau transformasi bidang \( \Re^2 \) berikut.
     \begin{equation*}
       \colvec{x \\ y} \mapsto \colvec{x/2 \\ y/3}
     \end{equation*}
     Tentukan citra elips berikut di bawah pemetaan ini.
     \begin{equation*}
       \set{\colvec{x \\ y} \suchthat (x^2/4)+(y^2/9)=1}
     \end{equation*}
     \begin{answer}
       Dengan \( \hat{x}=x/2 \) dan \( \hat{y}=y/3 \),
       himpunan citranya adalah
       \begin{equation*}
         \set{\colvec{\hat{x} \\ \hat{y}} \suchthat
           \frac{\displaystyle (2\hat{x})^2}{\displaystyle 4}
           +\frac{\displaystyle (3\hat{y})^2}{\displaystyle 9}=1}
         =\set{\colvec{\hat{x} \\ \hat{y}} \suchthat
            \hat{x}^2+\hat{y}^2=1}
       \end{equation*}
       yaitu lingkaran satuan pada bidang-\( \hat{x}\hat{y} \).
     \end{answer}
  \recommended \item
    Bayangkan seutas tali dililitkan rapat mengelilingi khatulistiwa bumi
    (anggap bumi berbentuk bola).
    Berapa banyak tali tambahan yang harus kita berikan agar, di sepanjang 
    keliling bumi, tali itu berada enam kaki di atas permukaan tanah?
    \begin{answer}
      Fungsi keliling $r\mapsto 2\pi r$ bersifat linear.  
      Dengan demikian, kita memperoleh
      $2\pi\cdot (r_{\text{earth}}+6)-
         2\pi\cdot (r_{\text{earth}})=12\pi$.
      Perhatikan bahwa
      jumlah tali tambahan yang diperlukan untuk menaikkan lingkaran dari lilitan rapat 
      pada bola basket hingga enam kaki di atas bola tersebut sama dengan jumlah yang diperlukan
      untuk menaikkannya dari lilitan rapat pada bumi hingga enam kaki di atas
      bumi.
     \end{answer}
  \recommended \item 
    Buktikan bahwa pemetaan \( \map{h}{\Re^3}{\Re} \) berikut
    \begin{equation*}
      \colvec{x \\ y \\ z}\;\mapsto\;
      \colvec{x \\ y \\ z}\dotprod\colvec[r]{3 \\ -1 \\ -1}=3x-y-z
    \end{equation*}
    bersifat linear.
    Perumum hasilnya.
    \begin{answer}
      Pembuktian bahwa pemetaan ini linear bersifat rutin.
      \begin{align*}
        h(c_1\cdot \colvec{x_1 \\ y_1 \\ z_1}
          +c_2\cdot \colvec{x_2 \\ y_2 \\ z_2})
        &=h(\colvec{c_1x_1+c_2x_2 \\ c_1y_1+c_2y_2 \\ c_1z_1+c_2z_2})   \\
        &=3(c_1x_1+c_2x_2)-(c_1y_1+c_2y_2)-(c_1z_1+c_2z_2)             \\
        &=c_1\cdot (3x_1-y_1-z_1)+c_2\cdot (3x_2-y_2-z_2)             \\
        &=c_1\cdot h(\colvec{x_1 \\ y_1 \\ z_1})
          +c_2\cdot h(\colvec{x_2 \\ y_2 \\ z_2})
      \end{align*}
      Dugaan perumuman yang wajar 
      adalah bahwa untuk sembarang \( \vec{k}\in\Re^3 \) yang tetap,
      pemetaan \( \vec{v}\mapsto\vec{v}\dotprod\vec{k} \) bersifat linear.
      Pernyataan ini benar.
      Hal ini mengikuti sifat-sifat hasil kali titik yang telah kita lihat sebelumnya:
      \( (\vec{v}+\vec{u})\dotprod\vec{k}=\vec{v}\dotprod\vec{k}+
         \vec{u}\dotprod\vec{k} \)
      dan \( (r\vec{v})\dotprod\vec{k}=r(\vec{v}\dotprod\vec{k}) \).
      (Dugaan wajar untuk memperumum perumuman ini, 
      yaitu bahwa pemetaan dari \( \Re^n \)
      ke \( \Re \) yang beraksi dengan mengambil hasil kali titik antara
      argumennya dan vektor tetap \( \vec{k}\in\Re^n \) bersifat linear, 
      juga benar.) 
    \end{answer}
  \item \label{exer:HomoRONeMultByScalar}
    Tunjukkan bahwa setiap homomorfisme dari \( \Re^1 \) ke \( \Re^1 \) 
    beraksi melalui perkalian dengan suatu skalar.
    Simpulkan bahwa setiap transformasi linear nontrivial pada \( \Re^1 \) merupakan
    isomorfisme.
    Apakah hal itu benar untuk transformasi pada \( \Re^2 \)?
    \( \Re^n \)?
    \begin{answer}
      Misalkan \( \map{h}{\Re^1}{\Re^1} \) linear.
      Suatu pemetaan linear ditentukan oleh aksinya pada basis, maka tetapkan basis
      \( \sequence{1} \) untuk \( \Re^1 \).
      Untuk sembarang \( r\in\Re^1 \), kita memiliki \( h(r)=h(r\cdot 1)=r\cdot h(1) \),
      sehingga \( h \) beraksi pada sembarang argumen $r$ 
      dengan mengalikannya dengan konstanta \( h(1) \).
      Jika \( h(1) \) tidak nol, maka pemetaan itu merupakan korespondensi\Dash inversnya 
      adalah pembagian dengan \( h(1) \)\Dash sehingga setiap transformasi nontrivial 
      pada $\Re^1$ merupakan isomorfisme.

      Pemetaan proyeksi berikut merupakan contoh yang menunjukkan bahwa tidak setiap
      transformasi pada \( \Re^n \) beraksi melalui perkalian dengan suatu konstanta
      ketika \( n>1 \), termasuk ketika $n=2$.
      \begin{equation*}
        \colvec{x_1 \\ x_2 \\ \vdots \\ x_n}
          \mapsto\colvec{x_1 \\ 0 \\ \vdots \\ 0}
      \end{equation*}  
    \end{answer}
  \item 
    %(This will be used in \nearbyexercise{exer:Cosets} below.)
    Tunjukkan bahwa untuk sembarang skalar \( a_{1,1},\dots, a_{m,n}  \), 
    pemetaan
    \( \map{h}{\Re^n}{\Re^m} \) berikut merupakan homomorfisme.
    \begin{equation*}
      \colvec{x_1 \\ \vdots \\ x_n}
        \mapsto
      \colvec{a_{1,1}x_1+\dots+a_{1,n}x_n \\ 
                   \vdots \\
                   a_{m,1}x_1+\cdots+a_{m,n}x_n}
    \end{equation*}
    \begin{answer}
      Dengan \( c \) dan \( d \) skalar, kita memperoleh hasil berikut.
          \begin{multline*}
            h(c\cdot \colvec{x_1 \\ \vdots \\ x_n} 
              +d\cdot \colvec{y_1 \\ \vdots \\ y_n})                     \\
            \begin{aligned}
            &=h(\colvec{cx_1+dy_1 \\ \vdots \\ cx_n+dy_n})        \\
            &=\colvec{a_{1,1}(cx_1+dy_1)+\dots+a_{1,n}(cx_n+dy_n) \\
                         \vdots                                   \\
                         a_{m,1}(cx_1+dy_1)+\dots+a_{m,n}(cx_n+dy_n)} \\
            &=c\cdot\colvec{a_{1,1}x_1+\dots+a_{1,n}x_n \\ 
                         \vdots                     \\ 
                         a_{m,1}x_1+\dots+a_{m,n}x_n}
            +d\cdot\colvec{a_{1,1}y_1+\dots+a_{1,n}y_n \\ 
                         \vdots \\ 
                         a_{m,1}y_1+\dots+a_{m,n}y_n} \\
            &=c\cdot h(\colvec{x_1 \\ \vdots \\ x_n})
              +d\cdot h(\colvec{y_1 \\ \vdots \\ y_n})
            \end{aligned}
          \end{multline*}
     \end{answer}
   \item Tinjau ruang polinomial \( \polyspace_n \).
     \begin{exparts}
       \partsitem Tunjukkan bahwa untuk setiap $i$, operator turunan ke-\( i \) 
          $d^i/dx^i$ merupakan transformasi linear pada ruang tersebut.
       \partsitem
          Simpulkan bahwa untuk sembarang skalar \( c_k,\ldots, c_0 \), pemetaan berikut
          merupakan transformasi linear pada ruang tersebut.
          \begin{equation*}
            f\;\mapsto\;
            c_{k}\frac{d^k}{dx^k}f+c_{k-1}\frac{d^{k-1}}{dx^{k-1}}f
               +\dots+
            c_1\frac{d}{dx}f+c_0f
          \end{equation*}
      \end{exparts}
      \begin{answer}
        \begin{exparts}
          \partsitem Setiap pangkat $i$ dari operator turunan bersifat linear 
            karena aturan-aturan berikut sudah dikenal dari kalkulus.
            \begin{equation*}
              \frac{d^i}{dx^i}(\,f(x)+g(x)\,)=\frac{d^i}{dx^i}f(x)
                                         +\frac{d^i}{dx^i}g(x)
              \qquad
              \frac{d^i}{dx^i}\,r\cdot f(x)=r\cdot\frac{d^i}{dx^i}f(x)
            \end{equation*}
         \partsitem 
          Setiap kombinasi linear dari pemetaan-pemetaan linear juga merupakan pemetaan linear.
          Dengan demikian,
          pemetaan yang diberikan merupakan transformasi linear pada \( \polyspace_n \).
        \end{exparts}
      \end{answer}
  \item 
    \nearbylemma{le:SpLinFcns} menunjukkan bahwa jumlah fungsi-fungsi linear bersifat
    linear dan bahwa kelipatan skalar suatu fungsi linear bersifat linear.
    Tunjukkan pula bahwa komposisi fungsi-fungsi linear bersifat linear.
    \begin{answer}
      (Argumen ini telah muncul sebagai bagian dari bukti bahwa
      isomorfisme merupakan relasi ekuivalensi.)
      Misalkan $\map{f}{U}{V}$ dan $\map{g}{V}{W}$ linear.
      Komposisinya mempertahankan kombinasi linear,
      \begin{multline*}
        \composed{g}{f}(c_1\vec{u}_1+c_2\vec{u}_2)
        =g(\,f(c_1\vec{u}_1+c_2\vec{u}_2)\,)          
        =g(\,c_1f(\vec{u}_1)+c_2f(\vec{u}_2)\,)              \\          
        =c_1\cdot g(f(\vec{u}_1))+c_2\cdot g(f(\vec{u}_2))          
        =c_1\cdot \composed{g}{f}(\vec{u}_1)
           +c_2\cdot \composed{g}{f}(\vec{u}_2)          
      \end{multline*}
      dengan $\vec{u}_1,\vec{u}_2\in U$ dan skalar $c_1,c_2$. 
    \end{answer}
  \item
    Misalkan \( \map{f}{V}{W} \) linear dan
    \( f(\vec{v}_1)=\vec{w}_1 \), \ldots, \( f(\vec{v}_n)=\vec{w}_n \)
    untuk sejumlah vektor \( \vec{w}_1 \), \ldots, \( \vec{w}_n \) dalam \( W \).
    \begin{exparts}
      \partsitem Jika himpunan vektor \( \vec{w} \) bebas, haruskah
        himpunan vektor \( \vec{v} \) juga bebas?
      \partsitem Jika himpunan vektor \( \vec{v} \) 
        bebas, haruskah
        himpunan vektor \( \vec{w} \) juga bebas?
      \partsitem Jika himpunan vektor \( \vec{w} \) merentang \( W \), haruskah himpunan
        vektor \( \vec{v} \) merentang \( V \)?
      \partsitem Jika himpunan vektor \( \vec{v} \) merentang \( V \), haruskah himpunan
        vektor \( \vec{w} \) merentang \( W \)?
    \end{exparts}
    \begin{answer}
      \begin{exparts}
        \partsitem Ya.
          Himpunan vektor $\vec{w}$ tidak mungkin bebas linear jika
          himpunan vektor $\vec{v}$ bergantung linear karena
          setiap relasi nontrivial dalam domain
          \( \zero_V=c_1\vec{v}_1+\dots+c_n\vec{v}_n \)
          akan menghasilkan relasi nontrivial dalam ruang hasil,
          \( f(\zero_V)=\zero_W=f(c_1\vec{v}_1+\dots+c_n\vec{v}_n)
             =c_1f(\vec{v}_1)+\dots+c_nf(\vec{v}_n)
             =c_1\vec{w}_1+\dots+c_n\vec{w}_n \).
        \partsitem Tidak selalu.
          Sebagai contoh, transformasi pada \( \Re^2 \) yang diberikan oleh
          \begin{equation*}
            \colvec{x \\ y} \mapsto \colvec{x+y \\ x+y}
          \end{equation*}
          memetakan himpunan bebas linear dalam domain berikut 
          ke citra yang bergantung linear.
          \begin{equation*}
            \set{\vec{v}_1,\vec{v}_2}=\set{\colvec[r]{1 \\ 0},\colvec[r]{1 \\ 1}}
            \;\mapsto\;
            \set{\colvec[r]{1 \\ 1},\colvec[r]{2 \\ 2}}=\set{\vec{w}_1,\vec{w}_2}
          \end{equation*}
        \partsitem Tidak selalu.
          Salah satu contohnya adalah pemetaan proyeksi 
          \( \map{\pi}{\Re^3}{\Re^2} \)
          \begin{equation*}
            \colvec{x \\ y \\ z}\mapsto\colvec{x \\ y}
          \end{equation*}
          dan himpunan berikut yang tidak merentang domain tetapi 
          dipetakan ke himpunan yang merentang kodomain.
          \begin{equation*}
            \set{\colvec[r]{1 \\ 0 \\ 0},\colvec[r]{0 \\ 1 \\ 0}}
            \mapsunder{\pi}\set{\colvec[r]{1 \\ 0},\colvec[r]{0 \\ 1}}
          \end{equation*}
        \partsitem Tidak selalu.
          Sebagai contoh, pemetaan inklusi $\map{\iota}{\Re^2}{\Re^3}$ memetakan
          basis standar $\stdbasis_2$ untuk domain ke himpunan yang 
          tidak merentang
          kodomain. 
          (\textit{Catatan.}
          Namun, himpunan vektor $\vec{w}$ memang merentang ruang hasil.
          Buktinya mudah.)
      \end{exparts}  
     \end{answer}
  \item  
    Perumum \nearbyexample{ex:MatTransMapLinear}
    dengan membuktikan bahwa untuk setiap domain dan kodomain yang sesuai,
    pemetaan transpos matriks bersifat linear.
    Apa domain dan kodomain yang sesuai itu?
    \begin{answer}
      Ingat bahwa entri pada baris~$i$ dan kolom~$j$ dari transpos $M$
      adalah entri $m_{j,i}$ dari baris~$j$ dan kolom~$i$ pada $M$.
      Sekarang, pemeriksaannya rutin.
      Mulailah dengan transpos kombinasi tersebut. 
      \begin{equation*}
        \trans{[r\cdot\begin{mat}
                \    &\vdots          \\
             \cdots &a_{i,j} &\cdots \\
                    &\vdots
           \end{mat}
         +s\cdot\begin{mat}
                \    &\vdots          \\
             \cdots &b_{i,j} &\cdots \\
                    &\vdots
           \end{mat}]} 
      \end{equation*}
      Gabungkan, lalu ambil transposnya.               
      \begin{equation*}
        =\trans{\begin{mat}
                 \       &\vdots                    \\
                  \cdots &ra_{i,j}+sb_{i,j} &\cdots \\
                         &\vdots
                \end{mat}}                               
        =\begin{mat}
              \     &\vdots                    \\
            \cdots &ra_{j,i}+sb_{j,i} &\cdots \\
                   &\vdots
          \end{mat}                               
      \end{equation*}
      Kemudian keluarkan skalar-skalar itu, lalu balikkan transposnya.
      \begin{multline*} 
        =r\cdot\begin{mat}
             \     &\vdots          \\
            \cdots &a_{j,i} &\cdots \\
                   &\vdots
          \end{mat}
         +s\cdot\begin{mat}
             \      &\vdots          \\
            \cdots &b_{j,i} &\cdots \\
                   &\vdots
          \end{mat}                             \\   
        =r\cdot\trans{\begin{mat}
             \      &\vdots          \\
            \cdots &a_{j,i} &\cdots \\
                   &\vdots
          \end{mat} }
         +s\cdot\trans{\begin{mat}
              \     &\vdots          \\
            \cdots &b_{j,i} &\cdots \\
                   &\vdots
          \end{mat} }
      \end{multline*}
      Domainnya adalah \( \matspace_{\nbym{m}{n}} \), 
      sedangkan kodomainnya adalah \( \matspace_{\nbym{n}{m}} \).  
     \end{answer}
  \item 
    \begin{exparts}
     \partsitem Untuk \( \vec{u},\vec{v}\in \Re^n \), 
        menurut definisi, ruas garis yang menghubungkan keduanya adalah himpunan
        \( \ell=\set{t\cdot\vec{u}+(1-t)\cdot\vec{v}\suchthat t\in [0..1]} \).
        Tunjukkan bahwa citra ruas antara
        $\vec{u}$ dan $\vec{v}$ di bawah homomorfisme $h$ adalah ruas antara $h(\vec{u})$ dan
        $h(\vec{v})$.
     \partsitem Himpunan bagian dari \( \Re^n \) disebut \definend{konveks}\index{himpunan konveks} 
       jika, untuk sembarang dua titik dalam
       himpunan tersebut, seluruh ruas garis yang menghubungkannya terletak dalam himpunan itu.
       (Bagian dalam suatu bola bersifat konveks,
       sedangkan kulit bolanya tidak.)
       Buktikan bahwa pemetaan linear dari \( \Re^n \) ke \( \Re^m \) mempertahankan 
       sifat konveks suatu himpunan.
     \end{exparts}
    \begin{answer}
      \begin{exparts}
        \partsitem  Untuk sembarang homomorfisme \( \map{h}{\Re^n}{\Re^m} \), kita memiliki
          \begin{equation*}
            h(\ell)
            =\set{h(t\cdot\vec{u}+(1-t)\cdot\vec{v})\suchthat t\in [0..1]}
            =\set{t\cdot h(\vec{u})+(1-t)\cdot h(\vec{v})\suchthat t\in [0..1]}
          \end{equation*}
          yang merupakan ruas garis dari $h(\vec{u})$ ke $h(\vec{v})$.
        \partsitem Kita harus menunjukkan bahwa jika suatu himpunan bagian dari domain bersifat konveks, maka
          citranya sebagai himpunan bagian dari ruang hasil juga konveks. 
          Misalkan \( C\subseteq \Re^n \) konveks,
          dan tinjau citranya $h(C)$.
          Untuk menunjukkan bahwa $h(C)$ konveks, kita harus menunjukkan bahwa untuk sembarang dua 
          anggotanya, $\vec{d}_1$ dan $\vec{d}_2$, 
          ruas garis yang menghubungkannya
          \begin{equation*}
            \ell=\set{t\cdot\vec{d}_1+(1-t)\cdot\vec{d}_2\suchthat t\in [0..1]}
          \end{equation*}
          merupakan himpunan bagian dari $h(C)$.

          Tetapkan sembarang anggota $\hat{t}\cdot\vec{d}_1+(1-\hat{t})\cdot\vec{d}_2$
          dari ruas garis tersebut.
          Karena titik-titik ujung $\ell$ berada dalam citra $C$, terdapat
          anggota-anggota $C$ yang dipetakan ke titik-titik tersebut, katakanlah $h(\vec{c}_1)=\vec{d}_1$
          dan $h(\vec{c}_2)=\vec{d}_2$.
          Sekarang, dengan $\hat{t}$ skalar yang kita tetapkan pada kalimat pertama
          paragraf ini, perhatikan bahwa
          $h(\hat{t}\cdot\vec{c}_1+(1-\hat{t})\cdot\vec{c}_2)
          =\hat{t}\cdot h(\vec{c}_1)+(1-\hat{t})\cdot h(\vec{c}_2)
          =\hat{t}\cdot\vec{d}_1+(1-\hat{t})\cdot\vec{d}_2$
          Dengan demikian, setiap anggota $\ell$ merupakan anggota $h(C)$, sehingga $h(C)$
          konveks.
      \end{exparts}
    \end{answer}
  \recommended \item \label{exer:HomosPresLinStruc}
    Misalkan \( \map{h}{\Re^n}{\Re^m} \) suatu homomorfisme.
    \begin{exparts}
      \partsitem Tunjukkan bahwa citra suatu garis dalam \( \Re^n \) 
        di bawah \( h \) adalah garis (yang mungkin degenerat) dalam \( \Re^m \).
      \partsitem Apa yang terjadi pada permukaan linear berdimensi \( k \)?
    \end{exparts}
    \begin{answer}
      \begin{exparts}
        \partsitem Untuk \( \vec{v}_0,\vec{v}_1\in\Re^n \), garis melalui 
          \( \vec{v}_0 \) dengan arah \( \vec{v}_1 \) adalah himpunan 
          $\set{\vec{v}_0+t\cdot \vec{v}_1\suchthat t\in\Re}$.
          Citra garis tersebut di bawah $h$,
          $\set{h(\vec{v}_0+t\cdot \vec{v}_1)\suchthat t\in\Re}
             =\set{h(\vec{v}_0)+t\cdot h(\vec{v}_1)\suchthat t\in\Re}$
          adalah garis melalui $h(\vec{v}_0)$ dengan arah $h(\vec{v}_1)$.
          Jika \( h(\vec{v}_1) \) adalah vektor nol, maka garis ini 
          degenerat.
        \partsitem Permukaan linear berdimensi \( k \) dalam \( \Re^n \) dipetakan ke
          permukaan linear berdimensi \( k \) dalam
          \( \Re^m \) (yang mungkin degenerat).
          Buktinya sama seperti bukti untuk garis.
      \end{exparts}  
     \end{answer}
  \item 
    Buktikan bahwa restriksi suatu homomorfisme ke subruang dari
    domainnya juga merupakan homomorfisme.
    \begin{answer}
      Misalkan \( \map{h}{V}{W} \) adalah homomorfisme dan misalkan
      \( S \) subruang dari \( V \).
      Tinjau pemetaan \( \map{\hat{h}}{S}{W} \) yang didefinisikan oleh
      \( \hat{h}(\vec{s})=h(\vec{s}) \).
      (Satu-satunya perbedaan antara $\hat{h}$ dan $h$ adalah perbedaan 
      domain.)
      Maka pemetaan baru ini linear: 
      \( \hat{h}(c_1\cdot\vec{s}_1+c_2\cdot\vec{s}_2)=
              h(c_1\vec{s}_1+c_2\vec{s}_2)=c_1h(\vec{s}_1)+c_2h(\vec{s}_2)=
              c_1\cdot\hat{h}(\vec{s}_1)+c_2\cdot\hat{h}(\vec{s}_2) \).  
    \end{answer}
  \item 
    Anggaplah \( \map{h}{V}{W} \) linear.
    \begin{exparts}
      \partsitem Tunjukkan bahwa \definend{ruang hasil} pemetaan ini, 
        \( \set{h(\vec{v})\suchthat \vec{v}\in V} \), merupakan subruang dari 
        kodomain \( W \).
      \partsitem Tunjukkan bahwa \definend{ruang nol} pemetaan ini,
        \( \set{\vec{v}\in V\suchthat h(\vec{v})=\zero_W} \) 
        merupakan subruang dari domain \( V \).
      \partsitem Tunjukkan bahwa jika \( U \) adalah subruang dari domain \( V \), maka
        citranya \( \set{h(\vec{u})\suchthat \vec{u}\in U} \) merupakan subruang 
        dari kodomain \( W \).
        Pernyataan ini memperumum butir pertama.
      \partsitem Perumum butir kedua.
    \end{exparts}
    \begin{answer} Hasil ini akan muncul sebagai lema dalam subbagian berikutnya.
      \begin{exparts}
        \partsitem Ruang hasil tidak kosong karena \( V \) tidak kosong.
          Untuk menuntaskan bukti, kita perlu menunjukkan bahwa ruang hasil tertutup terhadap kombinasi.
          Dengan \( \vec{v}_1,\dots,\vec{v}_n\in V \), suatu kombinasi vektor-vektor hasil
          berbentuk
          \begin{equation*}
            c_1\cdot h(\vec{v}_1)+\dots+c_n\cdot h(\vec{v}_n)
            =
            h(c_1\vec{v}_1)+\dots+h(c_n\vec{v}_n)
            =
            h(c_1\cdot \vec{v}_1+\dots+c_n\cdot \vec{v}_n),
          \end{equation*}
          yang berada dalam ruang hasil karena
          \( c_1\cdot \vec{v}_1+\dots+c_n\cdot \vec{v}_n \) merupakan anggota
          domain \( V \).
          Oleh karena itu, ruang hasil merupakan subruang.
        \partsitem Ruang nol tidak kosong karena memuat $\zero_V$, sebab
          \( \zero_V \) dipetakan ke \( \zero_W \).
          Ruang nol tertutup terhadap kombinasi linear karena, jika
          \( \vec{v}_1,\dots,\vec{v}_n\in V \) merupakan unsur-unsur 
          citra balik 
          \( \set{\vec{v}\in V\suchthat h(\vec{v})=\zero_W} \),
          maka untuk \( c_1,\ldots,c_n\in\Re \),
          \begin{equation*}
            \zero_W=c_1\cdot h(\vec{v}_1)+\dots+c_n\cdot h(\vec{v}_n)
            =h(c_1\cdot \vec{v}_1+\dots+c_n\cdot \vec{v}_n)
          \end{equation*}
          sehingga \( c_1\cdot \vec{v}_1+\dots+c_n\cdot \vec{v}_n \) juga berada dalam
          citra balik \( \zero_W \).
        \partsitem Citra \( U \) ini tidak kosong karena \( U \) tidak kosong.
          Untuk ketertutupan terhadap kombinasi, 
          dengan \( \vec{u}_1,\ldots,\vec{u}_n\in U \),
          \begin{equation*}
            c_1\cdot h(\vec{u}_1)+\dots+c_n\cdot h(\vec{u}_n)
            =
            h(c_1\cdot \vec{u}_1)+\dots+h(c_n\cdot \vec{u}_n)
            =
            h(c_1\cdot \vec{u}_1+\dots+c_n\cdot \vec{u}_n)
          \end{equation*}
          yang berada dalam \( h(U) \) karena
          \( c_1\cdot \vec{u}_1+\dots+c_n\cdot \vec{u}_n \) berada dalam \( U \).
          Jadi, himpunan ini merupakan subruang.
        \partsitem Perumuman yang wajar adalah bahwa citra balik suatu
          subruang kodomain merupakan subruang domain.

          Misalkan \( X \) adalah subruang dari \( W \).
          Perhatikan bahwa \( \zero_W\in X \), sehingga himpunan
          \( \set{\vec{v}\in V \suchthat h(\vec{v})\in X} \) tidak kosong.
          Untuk menunjukkan bahwa himpunan ini tertutup terhadap kombinasi, misalkan
          \( \vec{v}_1,\dots,\vec{v}_n \) merupakan unsur-unsur \( V \)
          sedemikian sehingga \( h(\vec{v}_1)=\vec{x}_1 \), \ldots,
          \( h(\vec{v}_n)=\vec{x}_n \), dan perhatikan bahwa 
          \begin{equation*}
            h(c_1\cdot \vec{v}_1+\dots+c_n\cdot \vec{v}_n)
            =c_1\cdot h(\vec{v}_1)+\dots+c_n\cdot h(\vec{v}_n)
            =c_1\cdot \vec{x}_1+\dots+c_n\cdot \vec{x}_n
          \end{equation*}
          sehingga kombinasi linear unsur-unsur \( h^{-1}(X) \) juga berada dalam
          \( h^{-1}(X) \).
      \end{exparts}  
     \end{answer}
  \item 
    Tinjau himpunan isomorfisme dari suatu ruang vektor ke dirinya sendiri.
    Apakah himpunan ini merupakan subruang dari ruang \( \linmaps{V}{V} \),
    yaitu ruang homomorfisme dari ruang tersebut ke dirinya sendiri?
    \begin{answer}
      Tidak; himpunan isomorfisme tidak memuat pemetaan nol
      (kecuali jika ruangnya trivial).
    \end{answer}
  \item 
   Apakah \nearbytheorem{th:HomoDetActOnBasis} mensyaratkan bahwa
   $\sequence{\vec{\beta}_1,\ldots,\vec{\beta}_n}$
   merupakan basis? 
   Dengan kata lain, apakah kita masih dapat memperoleh homomorfisme yang terdefinisi dengan baik dan tunggal jika kita
   membuang syarat bahwa himpunan vektor $\vec{\beta}$ 
   bebas linear, 
   atau syarat bahwa himpunan itu merentang domain?
   \begin{answer}
     Jika $\sequence{\vec{\beta}_1,\ldots,\vec{\beta}_n}$ tidak merentang ruang tersebut,
     maka pemetaannya tidak harus tunggal.
     Sebagai contoh, jika kita mencoba mendefinisikan pemetaan dari $\Re^2$ ke dirinya sendiri dengan 
     hanya menetapkan bahwa $\vec{e}_1$ dipetakan ke dirinya sendiri, maka 
     terdapat lebih dari
     satu homomorfisme yang mungkin; baik pemetaan identitas maupun pemetaan proyeksi 
     ke komponen pertama memenuhi syarat ini.

     Jika kita membuang syarat bahwa 
     $\sequence{\vec{\beta}_1,\ldots,\vec{\beta}_n}$
     bebas linear, maka kita berisiko memperoleh spesifikasi yang tidak konsisten
     (yaitu, pemetaan semacam itu mungkin tidak ada).
     Sebagai contoh, tinjau 
     $\sequence{\vec{e}_2,\vec{e}_1,2\vec{e}_1}$ dan cobalah 
     mendefinisikan pemetaan dari $\Re^2$ ke dirinya sendiri yang 
     memetakan $\vec{e}_2$ ke dirinya sendiri, serta memetakan
     $\vec{e}_1$ dan $2\vec{e}_1$ ke $\vec{e}_1$.
     Tidak ada homomorfisme yang dapat memenuhi ketiga syarat ini. 
   \end{answer}
  \item 
    Misalkan \( V \) suatu ruang vektor dan anggaplah bahwa 
    pemetaan \( \map{f_1,f_2}{V}{\Re^1} \) linear.
    \begin{exparts}
      \partsitem Definisikan pemetaan \( \map{F}{V}{\Re^2} \) yang fungsi-fungsi
        komponennya adalah fungsi linear yang diberikan.
        \begin{equation*}
           \vec{v}\mapsto\colvec{f_1(\vec{v}) \\ f_2(\vec{v})}
        \end{equation*}
        Tunjukkan bahwa \( F \) linear.
      \partsitem Apakah konversnya berlaku\Dash apakah setiap pemetaan linear dari \( V \) ke
        \( \Re^2 \) tersusun dari dua pemetaan komponen linear ke \( \Re^1 \)?
      \partsitem Perumum hasilnya.
    \end{exparts}
    \begin{answer}
      \begin{exparts}
        \partsitem Secara ringkas, pemeriksaan linearitasnya adalah sebagai berikut.
          \begin{multline*}
            F(r_1\cdot \vec{v}_1+r_2\cdot \vec{v}_2)
            =\colvec{f_1(r_1\vec{v}_1+r_2\vec{v}_2) \\ 
                        f_2(r_1\vec{v}_1+r_2\vec{v}_2)}           \\
            =r_1\colvec{f_1(\vec{v}_1) \\ f_2(\vec{v}_1)}
            +r_2\colvec{f_1(\vec{v}_2) \\ f_2(\vec{v}_2)}
            =r_1\cdot F(\vec{v}_1)+r_2\cdot F(\vec{v}_2)
          \end{multline*}
        \partsitem Ya.
          Misalkan \( \map{\pi_1}{\Re^2}{\Re^1} \) dan
          \( \map{\pi_2}{\Re^2}{\Re^1} \) adalah proyeksi
          \begin{equation*}
            \colvec{x \\ y}\mapsunder{\pi_1} x
              \quad\text{dan}\quad
            \colvec{x \\ y}\mapsunder{\pi_2} y
          \end{equation*}
          ke kedua sumbu.
          Sekarang, dengan \( f_1(\vec{v})=\pi_1(F(\vec{v})) \) dan
          \( f_2(\vec{v})=\pi_2(F(\vec{v})) \) 
          kita memperoleh fungsi-fungsi komponen yang diinginkan.
          \begin{equation*}
            F(\vec{v})=
            \colvec{f_1(\vec{v}) \\ f_2(\vec{v})}
          \end{equation*}
          Fungsi-fungsi itu linear karena merupakan komposisi fungsi-fungsi linear,
          dan fakta bahwa komposisi fungsi-fungsi linear bersifat linear
          merupakan bagian dari bukti bahwa isomorfisme adalah relasi
          ekuivalensi (sebagai alternatif, pemeriksaan bahwa fungsi-fungsi tersebut linear
          bersifat langsung). 
        \partsitem Secara umum, pemetaan dari ruang vektor \( V \) ke 
          \( \Re^n \) bersifat linear jika dan hanya jika setiap fungsi 
          komponennya linear.
          Pembuktiannya sama seperti pada butir sebelumnya.
      \end{exparts}  
     \end{answer}
\end{exercises}















\subsection{Ruang Hasil dan Ruang Nol}
Isomorfisme maupun homomorfisme sama-sama mempertahankan struktur.
Perbedaannya adalah bahwa homomorfisme memiliki lebih sedikit pembatasan karena
tidak harus pada dan
tidak harus satu-ke-satu.
Kita akan mengkaji apa yang dapat terjadi pada homomorfisme
tetapi tidak dapat terjadi pada isomorfisme. 

Pertama-tama, tinjau fakta bahwa 
homomorfisme tidak harus pada.
Tentu saja, 
setiap fungsi bersifat pada terhadap suatu himpunan, yaitu ruang hasilnya.
Sebagai contoh, pemetaan inklusi \( \map{\iota}{\Re^2}{\Re^3} \)
\begin{equation*}
  \colvec{x \\ y} \mapsto \colvec{x \\ y \\ 0}
\end{equation*}
merupakan homomorfisme dan tidak pada terhadap $\Re^3$. 
Namun, pemetaan ini pada terhadap bidang-$xy$.

\begin{lemma}   \label{le:RangeIsSubSp}
%<*lm:RangeIsSubSp>
Di bawah suatu homomorfisme, citra
setiap subruang dari domain merupakan subruang dari kodomain.
Secara khusus, citra seluruh ruang,
yaitu ruang hasil homomorfisme, merupakan subruang dari kodomain.
%</lm:RangeIsSubSp>
\end{lemma}

\begin{proof}
%<*pf:RangeIsSubSp>
Misalkan $\map{h}{V}{W}$ linear dan $S$ subruang dari domain 
$V$.
Citra $h(S)$ merupakan himpunan bagian dari kodomain $W$ yang
tidak kosong karena $S$ tidak kosong.
Jadi, untuk menunjukkan bahwa $h(S)$ merupakan subruang dari $W$,
kita hanya perlu menunjukkan bahwa himpunan itu tertutup terhadap kombinasi linear
dua vektor.
Jika $h(\vec{s}_1)$ dan $h(\vec{s}_2)$ merupakan anggota $h(S)$, maka
$c_1\cdot h(\vec{s}_1)+c_2\cdot h(\vec{s}_2)
  =
  h(c_1\cdot \vec{s}_1)+h(c_2\cdot \vec{s}_2)
  =
  h(c_1\cdot \vec{s}_1+c_2\cdot \vec{s}_2)$
juga merupakan anggota $h(S)$ karena merupakan citra dari
\( c_1\cdot \vec{s}_1+c_2\cdot \vec{s}_2 \) dalam \( S \).
%</pf:RangeIsSubSp>
\end{proof}

\begin{definition} \label{df:RangeSpace}
%<*df:RangeSpace>
\definend{Ruang hasil}\index{ruang hasil}\index{homomorfisme!ruang hasil}
dari suatu homomorfisme \( \map{h}{V}{W} \) adalah
\begin{equation*}
  \rangespace{h}=\set{h(\vec{v})\suchthat \vec{v}\in V}
\end{equation*}
yang kadang-kadang dinyatakan dengan \( h(V) \).
Dimensi ruang hasil disebut
\definend{rank} pemetaan tersebut.\index{rank!homomorfisme}
%</df:RangeSpace>
\end{definition}
\noindent 
Kita akan segera melihat hubungan antara rank suatu pemetaan dan rank suatu 
matriks.

\begin{example}  \label{ex:DerivMapRnge}
Untuk pemetaan turunan 
\( \map{d/dx}{\polyspace_3}{\polyspace_3} \)
yang diberikan oleh \( a_0+a_1x+a_2x^2+a_3x^3 \mapsto a_1+2a_2x+3a_3x^2 \),
ruang hasil 
\( \rangespace{d/dx} \) adalah himpunan polinomial kuadrat
\( \set{r+sx+tx^2\suchthat r,s,t\in\Re } \).
Dengan demikian, rank pemetaan ini adalah~\( 3 \).
\end{example}

\begin{example}  \label{ex:MatToPolyRnge}
Untuk homomorfisme \( \map{h}{M_{\nbyn{2}}}{\polyspace_3} \) berikut,
\begin{equation*}
  \begin{mat}
    a  &b  \\
    c  &d
  \end{mat}
  \mapsto
  (a+b+2d)+cx^2+cx^3
\end{equation*}
citra suatu vektor dalam ruang hasil
dapat memiliki sembarang suku konstan, harus memiliki koefisien $x$ sebesar nol,
dan harus memiliki koefisien $x^2$ yang sama dengan koefisien $x^3$.
Dengan kata lain, ruang hasilnya adalah
\( \rangespace{h}=\set{r+sx^2+sx^3\suchthat r,s\in\Re} \)
sehingga ranknya adalah~\( 2 \).
\end{example}

Hasil sebelumnya menunjukkan bahwa, 
dalam beralih dari definisi isomorfisme ke definisi homomorfisme yang
lebih umum,  
menghilangkan syarat pada tidak menimbulkan perbedaan yang mendasar.
Setiap homomorfisme bersifat pada terhadap suatu ruang, yaitu ruang hasilnya.

Namun, menghilangkan syarat satu-ke-satu memang menimbulkan perbedaan. 
Suatu homomorfisme dapat memiliki banyak unsur
domain yang dipetakan ke satu unsur kodomain.
Di bawah ini terdapat sketsa berbentuk kacang dari pemetaan 
banyak-ke-satu antardua himpunan.\appendrefs{pemetaan banyak-ke-satu}\spacefactor=1000 %
Sketsa tersebut menunjukkan tiga unsur kodomain yang masing-masing merupakan citra dari
banyak anggota domain.
(Alih-alih menggambar banyak panah $\mapsto$ satu per satu, setiap pengaitan
banyak masukan dengan satu keluaran hanya menampilkan satu panah semacam itu.)
\begin{center}  
  \includegraphics{map/mp/ch3.5}  % bean to bean; many to one
\end{center}
%<*InverseImage>
Ingat bahwa untuk sembarang fungsi $\map{h}{V}{W}$, 
himpunan unsur-unsur $V$ yang dipetakan ke \( \vec{w}\in W \)
disebut \definend{citra balik\/}\index{citra balik}%
\index{fungsi!citra balik}   
$h^{-1}(\vec{w})=\set{\vec{v}\in V\suchthat h(\vec{v})=\vec{w}}$.
%</InverseImage>
Di atas, sisi kiri menunjukkan tiga himpunan citra balik.

\begin{example}
Tinjau proyeksi\index{proyeksi}
\( \map{\pi}{\Re^3}{\Re^2} \)
\begin{equation*}
   \colvec{x \\ y \\ z}
    \mapsunder{\pi}
   \colvec{x \\ y}
\end{equation*}
yang merupakan homomorfisme banyak-ke-satu.
Suatu himpunan citra balik berupa garis vertikal yang terdiri dari vektor-vektor
dalam domain.
\begin{center}
  \includegraphics{map/mp/ch3.11}
\end{center}
Berikut salah satu contohnya.
\begin{equation*}
  \pi^{-1}(\colvec[r]{1 \\ 3})=\set{\colvec[r]{1 \\ 3 \\ z}\suchthat z\in\Re}
\end{equation*}
\end{example}

\begin{example} \label{ex:RTwoHomoREasyOneMap}
Homomorfisme $\map{h}{\Re^2}{\Re^1}$ berikut
\begin{equation*}
  \colvec{x \\ y}
   \mapsunder{h}
  x+y
\end{equation*}
juga bersifat banyak-ke-satu.
Untuk $w\in\Re^1$ yang tetap, 
citra balik $h^{-1}(w)$
\begin{center}
  \includegraphics{map/mp/ch3.12}
\end{center}
adalah himpunan vektor bidang yang jumlah komponennya sama dengan $w$.
\end{example}

% The above examples have only to do with the
% fact that we are considering functions,
% specifically, many-to-one functions.
% They show the inverse images 
% as sets of vectors that are 
% related to the image vector $\vec{w}$.
% But these are more than just arbitrary functions, they are
% homomorphisms; what do the two preservation
% conditions say about the relationships?

Ketika memperumum isomorfisme menjadi homomorfisme dengan 
menghilangkan syarat satu-ke-satu,
kita kehilangan sifat bahwa, secara intuitif, 
domain ``sama'' dengan ruang hasil.
Dengan kata lain, kita kehilangan korespondensi sempurna
antara domain dan ruang hasil.
Contoh-contoh di bawah ini menggambarkan bahwa hal yang tetap kita miliki
adalah bahwa suatu homomorfisme menjelaskan bagaimana
domain ``analog dengan'' atau ``serupa dengan'' ruang hasil.

\begin{example}    \label{ex:RThreeHomoRTwo} %\label{exPicProj}
Kita memandang $\Re^3$ serupa dengan $\Re^2$, kecuali bahwa vektor-vektornya memiliki satu
komponen tambahan. 
Dengan kata lain, kita memandang vektor berkomponen $x$, $y$, dan~$z$ 
serupa dengan vektor berkomponen $x$ dan~$y$.
Definisi pemetaan proyeksi $\pi$ memperjelas anggota
domain mana yang kita pandang berhubungan dengan anggota
kodomain mana.

Untuk memahami bagaimana
syarat-syarat pelestarian dalam definisi homomorfisme
menunjukkan bahwa unsur-unsur domain serupa dengan unsur-unsur kodomain, 
mulailah dengan membayangkan $\Re^2$ sebagai bidang-$xy$ di dalam $\Re^3$
(bidang-$xy$ di dalam $\Re^3$ 
adalah himpunan vektor tiga komponen dengan 
komponen ketiga
nol, sehingga tidak persis sama dengan himpunan vektor dua komponen~$\Re^2$,
tetapi penyematan ini membuat gambarnya jauh lebih jelas).
Sifat pelestarian penjumlahan menyatakan bahwa
vektor-vektor dalam \( \Re^3 \) berperilaku seperti bayangannya pada bidang tersebut.
\begin{center}  \small
  \begin{tabular}{@{}c@{}c@{}c@{}c@{}c@{}}
    \includegraphics{map/mp/ch3.1}
    &&\includegraphics{map/mp/ch3.2}
    &&\includegraphics{map/mp/ch3.3} \\[1.5ex]
    {\small  $\colvec{x_1 \\ y_1 \\ z_1}$ di atas $\colvec{x_1 \\ y_1}$}
    &{\small \ ditambah\ }
    &{\small $\colvec{x_2 \\ y_2 \\ z_2}$ di atas $\colvec{x_2 \\ y_2}$}
    &{\small \ sama dengan\ }
    &{\small $\colvec{x_1+x_2 \\ y_1+y_2 \\ z_1+z_2}$ 
              di atas $\colvec{x_1+x_2 \\ y_1+y_2}$}
  \end{tabular}
\end{center}
\noindent 
Memandang $\pi(\vec{v})$ sebagai ``bayangan'' $\vec{v}$ pada bidang 
memberikan pernyataan ulang berikut:
jumlah bayangan $\pi(\vec{v}_1)+\pi(\vec{v}_2)$ sama dengan
bayangan jumlahnya, 
$\pi(\vec{v}_1+\vec{v}_2)$.
Pelestarian perkalian skalar serupa.

Menggambar kodomain $\Re^2$ di sebelah kanan
menghasilkan gambar yang kurang menarik tetapi lebih sesuai dengan
sketsa berbentuk kacang di atas.
\begin{center}  \small
  \includegraphics{map/mp/ch3.4}
\end{center}
Sekali lagi, vektor-vektor domain yang dipetakan ke $\vec{w}_1$ terletak
pada suatu garis vertikal;
salah satunya digambar dengan warna abu-abu.
Sebut setiap anggota citra balik $\pi^{-1}(\vec{w}_1)$ ini
sebagai ``vektor~$\vec{w}_1$.''
Serupa dengan itu, terdapat garis vertikal ``vektor-vektor~$\vec{w}_2$'' dan
garis vertikal ``vektor-vektor~$\vec{w}_1+\vec{w}_2$.''
Sekarang, menyatakan bahwa $\pi$ merupakan homomorfisme berarti mengakui bahwa 
jika $\pi(\vec{v}_1)=\vec{w}_1$ dan $\pi(\vec{v}_2)=\vec{w}_2$, 
maka $\pi(\vec{v}_1+\vec{v}_2)=\pi(\vec{v}_1)+\pi(\vec{v}_2)
  =\vec{w}_1+\vec{w}_2$.
Dengan kata lain, kelas-kelas tersebut dapat dijumlahkan:~sembarang 
vektor~\( \vec{w}_1 \) ditambah sembarang
vektor~\( \vec{w}_2 \) 
menghasilkan vektor~\( \vec{w}_1+\vec{w}_2 \).
Perkalian skalarnya serupa.

Jadi, meskipun $\Re^3$ dan $\Re^2$ tidak isomorfik,
$\pi$ menjelaskan satu cara keduanya serupa:~vektor-vektor 
dalam $\Re^3$ dijumlahkan sebagaimana vektor-vektor terkait dalam $\Re^2$\Dash vektor-vektor 
dijumlahkan sebagaimana bayangannya dijumlahkan. 
\end{example}

\begin{example}   \label{ex:RTwoHomoRHardOne}
Suatu homomorfisme dapat menyatakan
analogi antarruang yang lebih halus daripada analogi sebelumnya.
Untuk pemetaan 
dari \nearbyexample{ex:RTwoHomoREasyOneMap},
\begin{equation*}
  \colvec{x \\ y}
   \mapsunder{h}
  x+y
\end{equation*}
tetapkan dua bilangan dalam ruang hasil, $w_1, w_2\in\Re$.
Suatu $\vec{v}_1$ yang dipetakan ke $w_1$ memiliki komponen-komponen yang 
berjumlah $w_1$,
sehingga citra balik $h^{-1}(w_1)$ adalah himpunan vektor 
yang titik ujungnya terletak pada garis diagonal $x+y=w_1$.
Pandang vektor-vektor ini sebagai ``vektor-vektor $w_1$.''
Serupa dengan itu, kita memiliki ``vektor-vektor $w_2$'' dan ``vektor-vektor $w_1+w_2$.''  
Sifat pelestarian penjumlahan menyatakan hal berikut.
\begin{center}  \small
  \begin{tabular}{@{}ccccc@{}}
    \includegraphics{map/mp/ch3.6}
    &&\includegraphics{map/mp/ch3.7}
    &&\includegraphics{map/mp/ch3.8}  \\[1.5ex]
    {\small sebuah ``vektor $w_1$''}
    &{\small ditambah}
    &{\small sebuah ``vektor $w_2$''}
    &{\small sama dengan}
    &{\small sebuah ``vektor $w_1+w_2$''}
  \end{tabular}
\end{center}
Dengan kata lain, jika kita menjumlahkan
vektor~$w_1$ dengan vektor~$w_2$, maka 
$h$ memetakan hasilnya ke vektor~$w_1+w_2$.
Secara ringkas, jumlah citra sama dengan citra jumlah.
Lebih ringkas lagi, \( h(\vec{v}_1)+h(\vec{v}_2)=h(\vec{v}_1+\vec{v}_2) \).
% (The preservation of scalar multiplication condition has a
% similar restatement.)
\end{example}

\begin{example}   \label{ex:PicRThreeToRTwo}
Citra balik dapat berupa struktur selain garis.
Untuk pemetaan linear \( \map{h}{\Re^3}{\Re^2} \),
\begin{equation*}
  \colvec{x \\ y \\ z}
    \mapsto
  \colvec{x \\ x}
\end{equation*}
himpunan-himpunan citra baliknya adalah bidang $x=0$, $x=1$, dan seterusnya,
yang tegak lurus terhadap sumbu-\( x \).
\begin{center}  \small
  \includegraphics{map/mp/ch3.9}
\end{center}
\end{example}

Kita tidak akan menjelaskan bagaimana setiap homomorfisme yang kita gunakan
merupakan suatu analogi karena makna formal
yang kita berikan kepada ``serupa dalam hal~\ldots'' adalah 
`terdapat suatu homomorfisme sedemikian sehingga~\ldots'.
Meskipun demikian, gagasan bahwa homomorfisme antardua ruang menjelaskan bagaimana
vektor-vektor domain terbagi ke dalam kelas-kelas yang berperilaku seperti
vektor-vektor ruang hasil merupakan cara yang baik untuk memandang homomorfisme.

Alasan lain kita tidak akan memperlakukan semua homomorfisme yang
kita jumpai seperti di atas adalah bahwa banyak ruang vektor sulit digambar,
misalnya ruang polinomial.
Namun, tidak ada salahnya memanfaatkan ruang-ruang
yang dapat kita gambar:
dari ketiga contoh
\ref{ex:RThreeHomoRTwo}, \ref{ex:RTwoHomoRHardOne}, 
dan~\ref{ex:PicRThreeToRTwo}, kita memperoleh dua wawasan.

Wawasan pertama adalah bahwa dalam ketiga contoh tersebut, 
citra balik vektor nol dalam ruang hasil merupakan garis atau bidang
yang melalui titik asal.
Karena itu, citra balik tersebut merupakan subruang dari domain.

\begin{lemma}  \label{le:NullspIsSubSp}
%<*lm:NullspIsSubSp>
Untuk sembarang homomorfisme, citra balik suatu subruang dari ruang hasil 
merupakan subruang dari domain.
Secara khusus, citra balik subruang trivial dari ruang hasil
merupakan subruang dari domain.
%</lm:NullspIsSubSp>
\end{lemma}

\noindent (Contoh-contoh di atas meninjau citra balik satu vektor,
tetapi hasil ini membahas citra balik himpunan, 
$h^{-1}(S)=\set{\vec{v}\in V\suchthat h(\vec{v})\in S}$.
Kita menggunakan istilah yang sama untuk keduanya dengan menganggap citra
balik suatu unsur tunggal $h^{-1}(\vec{w})$ sebagai
citra balik himpunan berunsur tunggal $h^{-1}(\set{\vec{w}})$.)

\begin{proof}
%<*pf:NullspIsSubSp>
Misalkan $\map{h}{V}{W}$ suatu homomorfisme
dan $S$ subruang dari ruang hasil $h$.
Tinjau citra balik $S$.
Himpunan ini tidak kosong karena memuat $\zero_V$, sebab
\( h(\zero_V)=\zero_W \) dan \( \zero_W \) merupakan unsur $S$
karena $S$ adalah subruang.
Untuk menuntaskan bukti, kita menunjukkan bahwa $h^{-1}(S)$ tertutup terhadap kombinasi linear.
Misalkan \( \vec{v}_1 \) dan \( \vec{v}_2 \) adalah dua unsurnya, sehingga
$h(\vec{v}_1)$ dan $h(\vec{v}_2)$ merupakan unsur-unsur $S$.
Maka
$c_1\vec{v}_1+c_2\vec{v}_2$
merupakan unsur citra balik $h^{-1}(S)$ 
karena
$h(c_1\vec{v}_1+c_2\vec{v}_2)
  =c_1h(\vec{v}_1)+c_2h(\vec{v}_2)$
merupakan anggota $S$. 
%</pf:NullspIsSubSp>
\end{proof}

\begin{definition} \label{df:NullSpace}
%<*df:NullSpace>
\definend{Ruang nol}\index{homomorfisme!ruang nol}\index{ruang nol}
atau \definend{kernel}\index{kernel, pemetaan linear} dari suatu pemetaan linear
\( \map{h}{V}{W} \) adalah citra balik $\zero_W$.
\begin{equation*}
  \nullspace{h}=h^{-1}(\zero_W)=\set{\vec{v}\in V\suchthat h(\vec{v})=\zero_W}
\end{equation*}
Dimensi ruang nol disebut
\definend{nulitas}\index{nulitas}\index{homomorfisme!nulitas} pemetaan tersebut.
%</df:NullSpace>
\end{definition}

\begin{center}
  \includegraphics{map/mp/ch3.10}
\end{center}

\begin{example}
Pemetaan dari \nearbyexample{ex:DerivMapRnge} memiliki ruang nol berikut,
\( \nullspace{d/dx}=\set{a_0+0x+0x^2+0x^3\suchthat a_0\in\Re} \)
sehingga nulitasnya adalah $1$.
\end{example}

\begin{example}
Pemetaan dari \nearbyexample{ex:MatToPolyRnge}
memiliki ruang nol berikut dan nulitas $2$.
\begin{equation*}
   \nullspace{h}=\set{\begin{mat}
                         a  &b          \\
                         0  &-(a+b)/2
                      \end{mat}\suchthat a,b\in\Re}
\end{equation*}
\end{example}

Sekarang kita beralih ke wawasan kedua dari contoh-contoh di atas.
Dalam \nearbyexample{ex:RThreeHomoRTwo}, setiap garis vertikal  
dipipihkan 
menjadi satu titik\Dash ketika beralih dari domain ke ruang hasil, $\pi$
mengambil semua garis vertikal satu dimensi ini dan memetakannya
ke satu titik,
sehingga dimensi ruang hasil lebih kecil satu daripada dimensi domain.
Serupa dengan itu, dalam \nearbyexample{ex:RTwoHomoRHardOne},
domain dua dimensi dimampatkan menjadi ruang hasil satu dimensi dengan membagi 
domain menjadi garis-garis diagonal 
dan memetakan masing-masing garis itu ke satu anggota ruang hasil.
Terakhir, dalam \nearbyexample{ex:PicRThreeToRTwo},
domain terbagi menjadi bidang-bidang yang
dipipihkan menjadi satu titik, sehingga pemetaan dimulai dengan domain tiga dimensi
tetapi berakhir dua dimensi lebih kecil, dengan
ruang hasil satu dimensi. 
(Kodomainnya
dua dimensi, tetapi ruang hasilnya satu dimensi, dan 
dimensi ruang hasillah yang penting.)

\begin{theorem} \label{th:RankPlusNullEqDim}
\index{rank!homomorfisme}
%<*th:RankPlusNullEqDim>
Rank suatu pemetaan linear ditambah nulitasnya sama dengan dimensi domainnya.
%</th:RankPlusNullEqDim>
\end{theorem}

\begin{proof}
%<*pf:RankPlusNullEqDim0>
Misalkan \( \map{h}{V}{W} \) linear dan
misalkan \( B_N=\sequence{\vec{\beta}_1,\ldots,\vec{\beta}_k} \)
adalah basis ruang nol.
Perluas basis itu menjadi basis
\( B_V=\sequence{\vec{\beta}_1,\dots,\vec{\beta}_k,
       \vec{\beta}_{k+1},\dots,\vec{\beta}_n} \)
untuk seluruh domain dengan menggunakan Korolari~Dua.III.\ref{cor:LIExpBas}.
Kita akan menunjukkan bahwa
\( B_R=\sequence{ h(\vec{\beta}_{k+1}),\dots,h(\vec{\beta}_n)} \)
merupakan basis ruang hasil.
Kemudian, menghitung ukuran basis-basis tersebut memberikan hasil yang diinginkan.
%</pf:RankPlusNullEqDim0>

%<*pf:RankPlusNullEqDim1>
Untuk melihat bahwa \( B_R \) bebas linear, 
tinjau 
\( \zero_W=c_{k+1}h(\vec{\beta}_{k+1})+\dots+c_nh(\vec{\beta}_n) \).
Kita memiliki
\( \zero_W=h(c_{k+1}\vec{\beta}_{k+1}+\dots+c_n\vec{\beta}_n) \)
sehingga \( c_{k+1}\vec{\beta}_{k+1}+\dots+c_n\vec{\beta}_n \)
berada dalam ruang nol $h$.
Karena \( B_N\) merupakan basis ruang nol, terdapat skalar-skalar
\( c_1,\dots,c_k \) yang memenuhi relasi berikut.
\begin{equation*}
   c_1\vec{\beta}_1+\dots+c_k\vec{\beta}_k
   =
   c_{k+1}\vec{\beta}_{k+1}+\dots+c_n\vec{\beta}_n
\end{equation*}
Namun, ini merupakan persamaan di antara anggota-anggota \( B_V \), 
yang merupakan basis \( V \), sehingga setiap $c_i$ sama dengan $0$.
Oleh karena itu, \( B_R \) bebas linear.
%</pf:RankPlusNullEqDim1>

%<*pf:RankPlusNullEqDim2>
Untuk menunjukkan bahwa \( B_R \) merentang ruang hasil,
tinjau suatu anggota ruang hasil \( h(\vec{v}) \).
Nyatakan \( \vec{v} \) sebagai kombinasi linear 
$\vec{v}=c_1\vec{\beta}_1+\dots+c_n\vec{\beta}_n$
dari anggota-anggota \( B_V \).
Hal ini memberikan
$h(\vec{v})=h(c_1\vec{\beta}_1+\dots+c_n\vec{\beta}_n)
     =c_1h(\vec{\beta}_1)+\dots+c_kh(\vec{\beta}_k)
     +c_{k+1}h(\vec{\beta}_{k+1})+\dots+c_nh(\vec{\beta}_n)$
dan karena 
$\vec{\beta}_1$, \ldots, $\vec{\beta}_k$ berada dalam ruang nol, kita memperoleh
$h(\vec{v})=\zero+\dots+\zero
     +c_{k+1}h(\vec{\beta}_{k+1})+\dots+c_nh(\vec{\beta}_n)$.
Dengan demikian, $h(\vec{v})$ merupakan kombinasi linear anggota-anggota \( B_R \),
sehingga $B_R$ merentang ruang hasil.
%</pf:RankPlusNullEqDim2>
\end{proof}

\begin{example}
Untuk \( \map{h}{\Re^3}{\Re^4} \) yang diberikan oleh
\begin{equation*}
  \colvec{x \\ y \\ z}
    \mapsunder{h}
  \colvec{x \\ 0 \\ y \\ 0}
\end{equation*}
ruang hasil dan ruang nolnya adalah
\begin{equation*}
  \rangespace{h}=
    \set{\colvec{a \\ 0 \\ b \\ 0}\suchthat a,b\in\Re }
  \quad\text{dan}\quad
  \nullspace{h}=
    \set{\colvec{0 \\ 0 \\ z}\suchthat z\in\Re }
\end{equation*}
sehingga rank $h$ adalah $2$, sedangkan nulitasnya adalah $1$.
\end{example}

\begin{example}
Jika \( \map{t}{\Re}{\Re} \) merupakan transformasi linear \( x\mapsto -4x, \)
maka ruang hasilnya adalah \(  \rangespace{t}=\Re \).
Ranknya adalah $1$ dan nulitasnya adalah $0$.
\end{example}

\begin{corollary}
\label{cor:RankDecreases}
Rank suatu pemetaan linear kurang dari atau sama dengan dimensi domainnya.
Kesamaan berlaku jika dan hanya jika nulitas pemetaan tersebut adalah $0$.
\end{corollary}

Kita mengetahui 
bahwa suatu isomorfisme terdapat di antara dua ruang 
jika dan hanya jika dimensi ruang hasil sama dengan dimensi domain.
Sekarang kita telah melihat bahwa syarat perlu bagi adanya homomorfisme adalah 
dimensi ruang hasil harus kurang dari atau sama dengan 
dimensi domain.
Sebagai contoh, tidak ada homomorfisme
dari \( \Re^2 \) yang pada terhadap \( \Re^3 \).
Ada banyak homomorfisme
dari \( \Re^2 \) ke dalam \( \Re^3 \), tetapi tidak ada yang pada.

Ruang hasil suatu pemetaan linear dapat memiliki dimensi yang lebih kecil secara ketat daripada 
dimensi domain,
sehingga
himpunan bebas linear dalam domain
dapat dipetakan ke himpunan bergantung linear dalam ruang hasil.
(Transformasi turunan pada $\polyspace_3$ dalam \nearbyexample{ex:DerivMapRnge}
memiliki domain berdimensi~$4$ tetapi ruang hasil berdimensi~$3$,
dan turunan memetakan
$\set{1,x,x^2,x^3}$ ke $\set{0,1,2x,3x^2}$).
Dengan kata lain, kebebasan dapat hilang di bawah homomorfisme.
Sebaliknya, kebergantungan tetap dipertahankan.

\begin{lemma} \label{lm:ImageLinearlyDependentIsLinearlyDependent}
%<*lm:ImageLinearlyDependentIsLinearlyDependent>
Di bawah suatu pemetaan linear, citra himpunan yang bergantung linear juga
bergantung linear.
%</lm:ImageLinearlyDependentIsLinearlyDependent>
\end{lemma}

\begin{proof}
%<*pf:ImageLinearlyDependentIsLinearlyDependent>
Misalkan \( c_1\vec{v}_1+\dots+c_n\vec{v}_n=\zero_V \)
dengan suatu $c_i$ tidak nol.
Terapkan $h$ pada kedua ruas:
\( h(c_1\vec{v}_1+\dots+c_n\vec{v}_n)=c_1h(\vec{v}_1)+\dots+c_nh(\vec{v}_n) \)
dan \( h(\zero_V)=\zero_W \).
Dengan demikian, kita memperoleh $c_1h(\vec{v}_1)+\dots+c_nh(\vec{v}_n)=\zero_W$ dengan suatu
$c_i$ tidak nol.
%</pf:ImageLinearlyDependentIsLinearlyDependent>
\end{proof}

Kapan kebebasan tidak hilang?
Syarat cukup yang jelas adalah ketika homomorfisme tersebut merupakan isomorfisme.
Syarat ini juga perlu; 
lihat \nearbyexercise{exer:NonSingIffPreservLI}.
Kita akan mengakhiri subbagian yang membandingkan homomorfisme dengan isomorfisme ini 
dengan mengamati bahwa
homomorfisme satu-ke-satu merupakan isomorfisme dari domainnya ke ruang hasilnya.

% \begin{definition}
% A linear map that is one-to-one is 
% {\em nonsingular}\index{homomorphism!nonsingular}%
% \index{nonsingular!homomorphism}.
% \end{definition}
% %\begin{remark}
% \noindent (In the next section we will see the connection
% between this use of `nonsingular' for maps and its familiar
% use for matrices.)

\begin{example}
Homomorfisme satu-ke-satu \( \map{\iota}{\Re^2}{\Re^3} \) berikut
\begin{equation*}
  \colvec{x \\ y}
    \mapsunder{\iota}
  \colvec{x \\ y \\ 0}
\end{equation*}
menghasilkan korespondensi di antara \( \Re^2 \) dan bidang-\( xy \),
yang merupakan himpunan bagian dari \( \Re^3 \).
\end{example}

% The prior observation allows us to adapt some results about isomorphisms.

\begin{theorem}  \label{th:OOHomoEquivalence}
%<*th:OOHomoEquivalence>
Jika \( V \) merupakan ruang vektor berdimensi \( n \), maka pernyataan-pernyataan
berikut ekuivalen
untuk pemetaan linear \( \map{h}{V}{W} \).
\begin{tfae}
  \item \( h \) bersifat satu-ke-satu
  \item \( h \) memiliki invers dari ruang hasilnya ke domainnya yang merupakan pemetaan 
         linear
  \item \( \nullspace{h}=\set{\zero\,} \), yaitu \( \nullity(h)=0 \)
  \item \( \rank (h)=n \)
  \item jika \( \sequence{\vec{\beta}_1,\dots,\vec{\beta}_n} \)
        merupakan basis untuk \( V \), maka
        \( \sequence{h(\vec{\beta}_1),\dots,h(\vec{\beta}_n)} \)
        merupakan basis untuk \( \rangespace{h} \)
\end{tfae}
%</th:OOHomoEquivalence>
\end{theorem}

\begin{proof}
%<*pf:OOHomoEquivalence0>
Pertama-tama kita akan menunjukkan bahwa \( \text{(1)} \Longleftrightarrow \text{(2)} \).
Kemudian kita akan menunjukkan bahwa
\(
  \text{(1)}\implies \text{(3)}\implies 
  \text{(4)}\implies \text{(5)}\implies \text{(2)}
\).
%</pf:OOHomoEquivalence0>

%<*pf:OOHomoEquivalence1>
Untuk \( \text{(1)} \Longrightarrow \text{(2)} \),
misalkan pemetaan linear $h$ bersifat satu-ke-satu sehingga memiliki invers
$\map{h^{-1}}{\rangespace{h}}{V}$.
Domain invers tersebut adalah ruang hasil $h$, sehingga kombinasi linear
dua anggotanya berbentuk $c_1h(\vec{v}_1)+c_2h(\vec{v}_2)$.
Pada kombinasi tersebut, invers \( h^{-1} \) memberikan hasil berikut.
\begin{align*}
  h^{-1}(c_1h(\vec{v}_1)+c_2h(\vec{v}_2))
  &=h^{-1}(h(c_1\vec{v}_1+c_2\vec{v}_2))  \\
  &=\composed{h^{-1}}{h}\;(c_1\vec{v}_1+c_2\vec{v}_2) \\
  &=c_1\vec{v}_1+c_2\vec{v}_2 \\
%   &=c_1\composed{h^{-1}}{h}\,(\vec{v}_1)
%            +c_2\composed{h^{-1}}{h}\,(\vec{v}_2) \\
  &=c_1\cdot h^{-1}(h(\vec{v}_1))+c_2\cdot h^{-1}(h(\vec{v}_2))
\end{align*}
Dengan demikian, jika suatu pemetaan linear memiliki invers, maka
invers tersebut harus linear.
Namun, hal ini juga memberikan implikasi \( \text{(2)} \Longrightarrow \text{(1)} \), 
karena invers itu sendiri harus satu-ke-satu.
%</pf:OOHomoEquivalence1>

%<*pf:OOHomoEquivalence2>
Dari implikasi-implikasi yang tersisa,
\( \text{(1)}\implies \text{(3)} \) berlaku karena setiap
homomorfisme memetakan \( \zero_V \) ke \( \zero_W \), tetapi pemetaan satu-ke-satu memetakan paling
banyak satu anggota \( V \) ke \( \zero_W \).

Selanjutnya, \( \text{(3)} \implies \text{(4)} \) benar karena rank
ditambah nulitas sama dengan dimensi domain.
%</pf:OOHomoEquivalence2>

%<*pf:OOHomoEquivalence3>
Untuk
\( \text{(4)} \implies \text{(5)} \), guna menunjukkan bahwa
\( \sequence{h(\vec{\beta}_1),\dots,h(\vec{\beta}_n)} \)
merupakan basis ruang hasil, kita hanya perlu menunjukkan bahwa barisan tersebut adalah himpunan perentang,
karena menurut asumsi ruang hasil berdimensi $n$.
Tinjau $h(\vec{v})\in\rangespace{h}$.
Menyatakan $\vec{v}$ sebagai kombinasi linear unsur-unsur basis menghasilkan
\( h(\vec{v})=h(\lincombo{c}{\vec{\beta}}) \),
yang memberikan 
\( h(\vec{v})=c_1h(\vec{\beta}_1)+\dots+c_nh(\vec{\beta}_n) \),
seperti yang diinginkan.
%</pf:OOHomoEquivalence3>

%<*pf:OOHomoEquivalence4>
Terakhir, untuk implikasi \( \text{(5)}\implies \text{(2)} \), anggaplah
\( \sequence{\vec{\beta}_1,\dots,\vec{\beta}_n} \)
merupakan basis untuk \( V \), sehingga
\( \sequence{h(\vec{\beta}_1),\dots,h(\vec{\beta}_n)} \)
merupakan basis untuk \( \rangespace{h} \).
Maka setiap
\( \vec{w}\in\rangespace{h} \) memiliki representasi tunggal
\( \vec{w}=c_1h(\vec{\beta}_1)+\dots+c_nh(\vec{\beta}_n) \).
Definisikan pemetaan dari \( \rangespace{h} \) ke $V$ dengan
\begin{equation*}
    \vec{w} \;\mapsto\; \lincombo{c}{\vec{\beta}}
\end{equation*}
(ketunggalan representasi menjadikan pemetaan ini terdefinisi dengan baik).
Pemeriksaan bahwa pemetaan ini linear dan bahwa
pemetaan ini merupakan invers $h$ bersifat mudah.
%</pf:OOHomoEquivalence3>
\end{proof}

Kita telah melihat bahwa suatu pemetaan linear
menyatakan bagaimana struktur domain serupa dengan struktur ruang hasil.
Kita dapat memandang pemetaan semacam itu sebagai pengorganisasian ruang domain menjadi
citra-citra balik dari titik-titik dalam ruang hasil.
Dalam kasus khusus ketika pemetaan tersebut satu-ke-satu, setiap citra balik berupa satu
titik dan pemetaan itu merupakan isomorfisme di antara domain dan ruang hasil.


\begin{exercises}
  \recommended \item 
    Misalkan \( \map{h}{\polyspace_3}{\polyspace_4} \) diberikan oleh
    \( p(x)\mapsto x\cdot p(x) \).
    Manakah di antara polinomial berikut yang berada dalam ruang nol?
    Manakah yang berada dalam ruang hasil?
    \begin{exparts*}
      \partsitem \( x^3 \)
      \partsitem \( 0 \)
      \partsitem \( 7 \)
      \partsitem \( 12x-0.5x^3 \)
      \partsitem \( 1+3x^2-x^3 \)
    \end{exparts*}
    \begin{answer}
      Pertama, untuk menjawab apakah suatu polinomial berada dalam ruang nol, 
      kita harus meninjaunya sebagai anggota domain $\polyspace_3$.
      Untuk menjawab apakah polinomial itu berada dalam ruang hasil, kita meninjaunya sebagai anggota
      kodomain $\polyspace_4$.
      Dengan kata lain, untuk $p(x)=x^4$, pertanyaan apakah polinomial itu berada dalam ruang hasil
      masuk akal, tetapi pertanyaan apakah polinomial itu berada dalam ruang nol tidak,
      karena polinomial tersebut bahkan tidak berada dalam domain.
      \begin{exparts}
        \partsitem Polinomial $x^3\in\polyspace_3$ tidak berada dalam ruang nol
            karena $h(x^3)=x^4$ bukan polinomial nol dalam $\polyspace_4$.
            Polinomial $x^3\in\polyspace_4$ berada dalam ruang hasil karena
            $x^2\in\polyspace_3$ dipetakan oleh $h$ ke $x^3$.
        \partsitem Jawaban untuk kedua pertanyaan adalah, ``Ya, karena 
            $h(0)=0$.''
            Polinomial $0\in\polyspace_3$ berada dalam ruang nol
            karena dipetakan oleh $h$ ke polinomial nol dalam 
            $\polyspace_4$.
            Polinomial $0\in\polyspace_4$ berada dalam ruang hasil karena
            merupakan citra $0\in\polyspace_3$ di bawah $h$.
        \partsitem Polinomial $7\in\polyspace_3$ tidak berada dalam ruang nol
            karena $h(7)=7x$ bukan polinomial nol dalam $\polyspace_4$.
            Polinomial $7\in\polyspace_4$ tidak berada dalam ruang hasil 
            karena tidak ada anggota domain yang, ketika dikalikan 
            dengan $x$, menghasilkan polinomial konstan $p(x)=7$.
        \partsitem Polinomial $12x-0.5x^3\in\polyspace_3$ tidak berada dalam 
            ruang nol karena $h(12x-0.5x^3)=12x^2-0.5x^4$.
            Polinomial $12x-0.5x^3\in\polyspace_4$ berada dalam ruang hasil 
            karena merupakan citra $12-0.5x^2$.
        \partsitem Polinomial $1+3x^2-x^3\in\polyspace_3$ tidak berada dalam 
            ruang nol karena $h(1+3x^2-x^3)=x+3x^3-x^4$.
            Polinomial $1+3x^2-x^3\in\polyspace_4$ tidak berada dalam 
            ruang hasil karena suku konstannya.
      \end{exparts}
    \end{answer}
  \item Tentukan ruang hasil dan rank setiap homomorfisme berikut.
    \begin{exparts}
      \partsitem
        $\map{h}{\polyspace_3}{\Re^2}$ yang diberikan oleh
        \begin{equation*}
           ax^2+bx+c\mapsto \colvec{a+b \\ a+c}
        \end{equation*}
      \partsitem $\map{f}{\Re^2}{\Re^3}$ yang diberikan oleh 
        \begin{equation*}
          \colvec{x \\ y}\mapsto\colvec{0 \\ x-y \\ 3y} 
        \end{equation*}     
    \end{exparts}
    \begin{answer}
      \begin{exparts}
        \partsitem
          Ruang hasil $h$ adalah seluruh kodomain~$\Re^2$ karena, untuk sembarang 
          \begin{equation*}
            \colvec{x \\ y}\in\Re^2
          \end{equation*}
          vektor itu merupakan citra dari vektor domain $0x^2+bx+c$ di bawah~$h$.
          Jadi, rank~$h$ adalah~$2$.
        \partsitem
          Ruang hasilnya adalah bidang-$yz$. Sembarang 
          \begin{equation*}
            \colvec{0 \\ a \\ b}
          \end{equation*}
          merupakan citra vektor domain berikut di bawah~$f$.
          \begin{equation*}
            \colvec{a+b/3  \\ b/3}
          \end{equation*}
          Jadi, rank pemetaan tersebut adalah~$2$.
      \end{exparts}
    \end{answer}
  \recommended \item
    Tentukan ruang hasil dan rank setiap pemetaan berikut.
    \begin{exparts}
      \partsitem \( \map{h}{\Re^2}{\polyspace_3} \) yang diberikan oleh
        \begin{equation*}
          \colvec{a \\ b}\mapsto a+ax+ax^2
        \end{equation*}
      \partsitem \( \map{h}{\matspace_{\nbyn{2}}}{\Re} \) yang diberikan oleh
        \begin{equation*}
          \begin{mat}
            a  &b  \\
            c  &d
          \end{mat}
          \mapsto a+d
        \end{equation*}
      \partsitem \( \map{h}{\matspace_{\nbyn{2}}}{\polyspace_2} \) yang diberikan oleh
        \begin{equation*}
          \begin{mat}
            a  &b  \\
            c  &d
          \end{mat}
          \mapsto a+b+c+dx^2
        \end{equation*}
      \partsitem pemetaan nol \( \map{Z}{\Re^3}{\Re^4} \)
    \end{exparts}
    \begin{answer}
      \begin{exparts}
        \partsitem Ruang hasilnya adalah
          \begin{equation*}
            \rangespace{h}
            =\set{a+ax+ax^2\in\polyspace_3\suchthat a\in\Re}
            =\set{a\cdot(1+x+x^2)\suchthat a\in\Re}
          \end{equation*}
          sehingga ranknya adalah satu.
        \partsitem Ruang hasil
          \begin{equation*}
            \rangespace{h}=
             \set{a+d\suchthat a,b,c,d\in\Re}
          \end{equation*}
          adalah seluruh $\Re$ (kita dapat memperoleh sembarang bilangan real dengan
mengambil $d=0$ dan $a$ sama dengan bilangan yang diinginkan).
          Dengan demikian, ranknya adalah satu.
        \partsitem Ruang hasilnya adalah 
          $\rangespace{h}=\set{r+sx^2\suchthat r,s\in\Re}$.
          Ranknya adalah dua.
        \partsitem Ruang hasilnya adalah subruang trivial dari \( \Re^4 \), 
          sehingga ranknya adalah nol.
      \end{exparts}  
    \end{answer}
  \recommended \item
    Untuk setiap pemetaan linear dalam latihan sebelumnya, tentukan ruang nol dan nulitasnya.
    \begin{answer}
      \begin{exparts}
        \partsitem Ruang nolnya adalah
          \begin{align*}
            \nullspace{h}
              &=\set{\colvec{a \\ b}\in\Re^2\suchthat 
                                    a+ax+ax^2+0x^3=0+0x+0x^2+0x^3}  \\
              &=\set{\colvec{0 \\ b}\suchthat b\in\Re}
          \end{align*}
        sehingga nulitasnya adalah satu.
        \partsitem Ruang nolnya adalah sebagai berikut.
          \begin{equation*}
            \nullspace{h}
             =\set{\begin{mat}
                     a  &b  \\
                     c  &d
                   \end{mat} \suchthat a+d=0}
             =\set{\begin{mat}
                     -d  &b  \\
                      c  &d
                   \end{mat} \suchthat b,c,d\in\Re }
          \end{equation*}
          Dengan demikian, nulitasnya adalah tiga.
        \partsitem Ruang nol
          \begin{equation*}
            \nullspace{h}
             =\set{\begin{mat}
                     a  &b  \\
                     c  &d
                   \end{mat} \suchthat a+b+c=0,\,d=0}
             =\set{\begin{mat}
                     -b-c  &b  \\
                      c  &0
                   \end{mat} \suchthat b,c\in\Re }
          \end{equation*}
          merupakan ruang berdimensi~$2$, sehingga nulitasnya adalah dua.
        \partsitem Setiap vektor dalam domain dipetakan ke vektor
          nol, sehingga ruang nolnya adalah $\nullspace{Z}=\Re^3$.
      \end{exparts}  
    \end{answer}
  \recommended \item 
    Tentukan nulitas setiap pemetaan di bawah ini.
    \begin{exparts*}
      \partsitem \( \map{h}{\Re^5}{\Re^8} \) dengan rank lima
      \partsitem \( \map{h}{\polyspace_3}{\polyspace_3} \) dengan rank satu
      \partsitem \( \map{h}{\Re^6}{\Re^3} \), suatu pemetaan pada
      \partsitem \( \map{h}{\matspace_{\nbyn{3}}}{\matspace_{\nbyn{3}}} \),
        pada
    \end{exparts*}
    \begin{answer}
      Untuk masing-masing, gunakan hasil 
      bahwa rank ditambah nulitas sama dengan dimensi domain.
      \begin{exparts*}
        \partsitem $0$
        \partsitem $3$ 
        \partsitem $3$
        \partsitem $0$
      \end{exparts*}
    \end{answer}
  \recommended \item 
    Apa ruang nol transformasi diferensiasi
    \( \map{d/dx}{\polyspace_n}{\polyspace_n} \)?
    Apa ruang nol turunan kedua sebagai
    transformasi pada \( \polyspace_n \)?
    Bagaimana dengan turunan ke-\( k \)?
    \begin{answer}
      Karena
      \begin{equation*}
        \frac{d}{dx}\,(a_0+a_1x+\cdots+a_nx^n)
        =a_1+2a_2x+3a_3x^2+\cdots+na_nx^{n-1}
      \end{equation*}
      kita memperoleh hasil berikut.
      \begin{align*}
        \nullspace{\frac{d}{dx}}
        &=\set{a_0+\cdots+a_nx^n\suchthat a_1+2a_2x+\cdots+na_nx^{n-1}
                                         =0+0x+\cdots+0x^{n-1}}       \\
        &=\set{a_0+\cdots+a_nx^n\suchthat 
               \text{$a_1=0$, dan $a_2=0$, \ldots, $a_n=0$}}  \\
        &=\set{a_0+0x+0x^2+\cdots+0x^n\suchthat a_0\in\Re}
      \end{align*}
      Dengan cara yang sama, 
      \begin{equation*}
        \nullspace{\frac{d^k}{dx^k}}
        =\set{a_0+a_1x+\cdots+a_{k-1}x^{k-1}\suchthat a_0,\ldots,a_{k-1}\in\Re}
      \end{equation*}
      untuk \( k\leq n \).  
    \end{answer}
  \item Untuk pemetaan $\map{h}{\Re^3}{\Re^2}$ yang diberikan oleh
      \begin{equation*}
        \colvec{x \\  y \\ z}
          \mapsto
          \colvec{x+y \\ x+z}
      \end{equation*}
      tentukan ruang hasil, rank, ruang nol, dan nulitasnya.
      \begin{answer}
        Untuk melihat bahwa ruang hasilnya adalah seluruh~$\Re^2$, perhatikan bahwa untuk sembarang
        $u,v\in\Re$, sistem berikut memiliki solusi.
        \begin{equation*}
          \begin{linsys}{3}
            x &+ &y &  &  &= &u  \\
            x &  &  &+ &z &= &v
          \end{linsys}
          \grstep{-\rho_1+\rho_2}
          \begin{linsys}{3}
            x &+ &y  &  &  &= &u  \\
              &  &-y &+ &z &= &-u+v
          \end{linsys}
        \end{equation*}
        (Bahkan, karena terdapat variabel bebas $z$, sistem ini memiliki tak berhingga 
        banyak solusi.)
        Dengan demikian, ranknya, yaitu dimensi ruang hasil, adalah~$2$.

        Karena rank ditambah nulitas sama dengan dimensi domain,
        kita mengetahui bahwa nulitasnya adalah~$1$.
        Menentukan ruang nol membuktikan hal tersebut:
        \begin{equation*}
          \begin{linsys}{3}
            x &+ &y &  &  &= &0  \\
            x &  &  &+ &z &= &0
          \end{linsys}
          \grstep{-\rho_1+\rho_2}
          \begin{linsys}{3}
            x &+ &y  &  &  &= &0  \\
              &  &-y &+ &z &= &0
          \end{linsys}
        \end{equation*}
        ruang nolnya adalah himpunan berikut,
        \begin{equation*}
          \set{\colvec{-1 \\ 1 \\ 1}\cdot z\suchthat z \in \Re}
        \end{equation*}
        yang berdimensi satu.
      \end{answer}
  \item  \label{exer:CondTwoProjMap}
    \nearbyexample{ex:RThreeHomoRTwo} menyatakan ulang syarat pertama dalam
    definisi homomorfisme sebagai `bayangan suatu jumlah adalah jumlah
    bayangan-bayangannya'.
    Nyatakan ulang syarat kedua dengan gaya yang sama.
    \begin{answer}
      Bayangan suatu kelipatan skalar adalah kelipatan skalar dari bayangannya.  
    \end{answer}
  \item 
    Untuk homomorfisme \( \map{h}{\polyspace_3}{\polyspace_3} \)
    yang diberikan oleh \( h(a_0+a_1x+a_2x^2+a_3x^3)=a_0+(a_0+a_1)x+(a_2+a_3)x^3 \),
    tentukan hal-hal berikut.
    \begin{exparts*}
      \partsitem \( \nullspace{h} \)
      \partsitem \( h^{-1}(2-x^3) \)
      \partsitem \( h^{-1}(1+x^2) \)
    \end{exparts*}
    \begin{answer}
      \begin{exparts}
        \partsitem Menetapkan $a_0+(a_0+a_1)x+(a_2+a_3)x^3=0+0x+0x^2+0x^3$
           memberikan $a_0=0$, $a_0+a_1=0$, dan $a_2+a_3=0$, sehingga ruang nolnya adalah
           \( \set{-a_3x^2+a_3x^3\suchthat a_3\in\Re} \).
        \partsitem Menetapkan $a_0+(a_0+a_1)x+(a_2+a_3)x^3=2+0x+0x^2-x^3$
           memberikan $a_0=2$, $a_1=-2$, dan $a_2+a_3=-1$.
           Mengambil $a_3$ sebagai parameter dan menamainya ulang $a_3=a$ memberikan
           deskripsi himpunan berikut,
           \( \set{2-2x+(-1-a)x^2+ax^3\suchthat a\in\Re}
                =\set{(2-2x-x^2)+a\cdot(-x^2+x^3)\suchthat a\in\Re} \).
        \partsitem Himpunan ini kosong karena ruang hasil $h$ hanya memuat
           polinomial-polinomial dengan suku $0x^2$.
      \end{exparts}  
    \end{answer}
  \recommended \item 
    Untuk pemetaan \( \map{f}{\Re^2}{\Re} \) yang diberikan oleh
    \begin{equation*}
       f(\colvec{x \\ y})=2x+y
    \end{equation*}
    gambarkan sketsa himpunan-himpunan citra balik berikut:~\( f^{-1}(-3) \), \( f^{-1}(0) \), 
    dan \( f^{-1}(1) \).
    \begin{answer}
      Semua citra balik merupakan garis dengan gradien $-2$.    
      \begin{center}
        \includegraphics{map/mp/ch3.74}
     \end{center}
    \end{answer}
  \recommended \item 
    Setiap transformasi pada \( \polyspace_3 \) berikut bersifat satu-ke-satu.
    Untuk masing-masing, tentukan inversnya.
    \begin{exparts}
      \partsitem \( a_0+a_1x+a_2x^2+a_3x^3\mapsto a_0+a_1x+2a_2x^2+3a_3x^3 \)
      \partsitem \( a_0+a_1x+a_2x^2+a_3x^3\mapsto a_0+a_2x+a_1x^2+a_3x^3 \)
      \partsitem \( a_0+a_1x+a_2x^2+a_3x^3\mapsto a_1+a_2x+a_3x^2+a_0x^3 \)
      \partsitem \( a_0+a_1x+a_2x^2+a_3x^3\mapsto
              a_0+(a_0+a_1)x+(a_0+a_1+a_2)x^2+(a_0+a_1+a_2+a_3)x^3 \)
    \end{exparts}
    \begin{answer}
      Berikut invers-inversnya.
      \begin{exparts}
        \partsitem \( a_0+a_1x+a_2x^2+a_3x^3\mapsto
                        a_0+a_1x+(a_2/2)x^2+(a_3/3)x^3 \)
        \partsitem \( a_0+a_1x+a_2x^2+a_3x^3\mapsto a_0+a_2x+a_1x^2+a_3x^3 \)
        \partsitem \( a_0+a_1x+a_2x^2+a_3x^3\mapsto a_3+a_0x+a_1x^2+a_2x^3 \)
        \partsitem \( a_0+a_1x+a_2x^2+a_3x^3\mapsto
                        a_0+(a_1-a_0)x+(a_2-a_1)x^2+(a_3-a_2)x^3 \)
      \end{exparts} 
      Sebagai contoh, untuk pemetaan kedua, pemetaan yang diberikan dalam soal memetakan
      $0+1x+2x^2+3x^3\mapsto 0+2x+1x^2+3x^3$, lalu invers di atas
      memetakan $0+2x+1x^2+3x^3\mapsto 0+1x+2x^2+3x^3$.
      Jadi, pemetaan ini sebenarnya merupakan invers bagi dirinya sendiri.
     \end{answer}
  \item 
    Deskripsikan ruang nol dan ruang hasil dari transformasi
    yang diberikan oleh \( \vec{v}\mapsto 2\vec{v} \).
    \begin{answer}
       Untuk sembarang ruang vektor $V$, ruang nol 
       \begin{equation*}
         \set{\vec{v}\in V\suchthat 2\vec{v}=\zero}
       \end{equation*}
       bersifat trivial, sedangkan ruang hasil 
       \begin{equation*}
         \set{\vec{w}\in V\suchthat 
                 \text{$\vec{w}=2\vec{v}$ untuk suatu $\vec{v}\in V$}}
       \end{equation*}
       adalah seluruh \( V \), karena setiap vektor $\vec{w}$ sama dengan dua kali suatu 
       vektor lain; secara khusus, vektor itu sama dengan dua kali $(1/2)\vec{w}$.
       (Dengan demikian, transformasi ini sebenarnya merupakan automorfisme.)  
    \end{answer}
  \item 
    Daftarkan semua pasangan \( (\text{rank}(h),\text{nulitas}(h)) \) 
    yang mungkin untuk pemetaan linear dari $\Re^5$ ke $\Re^3$.
    \begin{answer}
      Karena rank ditambah nulitas sama dengan dimensi domain 
      (di sini lima), dan rank paling tinggi tiga,
      pasangan yang mungkin adalah:~\( (3,2) \), \( (2,3) \), 
      \( (1,4) \), dan~\( (0,5) \).
      Mudah untuk menyusun pemetaan linear yang menunjukkan bahwa setiap pasangan tersebut memang
mungkin.  
    \end{answer}
  \item 
    Apakah pemetaan diferensiasi
    \( \map{d/dx}{\polyspace_n}{\polyspace_n} \) memiliki invers?
    \begin{answer}
      Tidak (kecuali \( \polyspace_n \) trivial), karena kedua polinomial
      \( f_0(x)=0 \) dan \( f_1(x)=1 \) memiliki turunan yang sama; suatu pemetaan harus
      satu-ke-satu agar memiliki invers.
    \end{answer}
  \recommended \item
    Tentukan nulitas pemetaan \( \map{h}{\polyspace_n}{\Re} \) berikut.
    \begin{equation*}
      a_0+a_1x+\dots+a_nx^n\mapsto\int_{x=0}^{x=1}a_0+a_1x+\dots+a_nx^n\,dx
    \end{equation*}
    \begin{answer}
      Ruang nolnya adalah sebagai berikut.
      \begin{multline*}
        \set{a_0+a_1x+\dots+a_nx^n\suchthat
                   a_0(1)+\frac{\displaystyle a_1}{\displaystyle 2}(1^2)
                   +\dots
                   +\frac{\displaystyle a_n}{\displaystyle n+1}(1^{n+1})=0} \\
        =\set{a_0+a_1x+\dots+a_nx^n\suchthat
                   a_0+(a_1/2)+\dots+(a_{n}/(n+1))=0}
      \end{multline*}
      Dengan demikian, nulitasnya adalah \( n \).  
     \end{answer}
  \item 
    \begin{exparts}
      \partsitem Buktikan bahwa suatu homomorfisme bersifat pada jika dan hanya jika
        ranknya sama dengan dimensi kodomainnya.
      \partsitem Simpulkan bahwa homomorfisme antarruang vektor dengan 
        dimensi yang sama bersifat satu-ke-satu jika dan hanya jika bersifat pada.
    \end{exparts}
    \begin{answer}
      \begin{exparts}
        \partsitem Satu arah sudah jelas:~jika homomorfisme tersebut 
          pada, maka ruang hasilnya adalah kodomain, sehingga
          ranknya sama dengan dimensi kodomainnya.
          Untuk arah lainnya, anggaplah rank pemetaan sama dengan
          dimensi kodomain.
          Maka ruang hasil pemetaan merupakan subruang dari kodomain dan memiliki
          dimensi yang sama dengan dimensi kodomain.
          Oleh karena itu, ruang hasil pemetaan harus sama dengan kodomain dan
          pemetaan tersebut bersifat pada. 
          (Kata `oleh karena itu' berlaku karena
          terdapat himpunan bagian bebas linear dari ruang hasil yang
          ukurannya sama dengan dimensi kodomain,
          tetapi setiap himpunan bagian bebas linear semacam itu dalam kodomain harus merupakan 
          basis kodomain, sehingga ruang hasil sama dengan kodomain.)
        \partsitem By \nearbytheorem{th:OOHomoEquivalence}, 
          suatu homomorfisme bersifat satu-ke-satu jika dan hanya jika nulitasnya nol.
          Karena rank ditambah nulitas sama dengan dimensi domain,
          dapat disimpulkan bahwa suatu homomorfisme bersifat satu-ke-satu jika dan hanya
          jika ranknya sama dengan dimensi domainnya.
          Namun, domain dan kodomain ini memiliki dimensi yang sama,
          sehingga pemetaan tersebut satu-ke-satu jika dan hanya jika pada.  
      \end{exparts}
   \end{answer}
  \item \label{exer:NonSingIffPreservLI} 
    Tunjukkan bahwa suatu pemetaan linear bersifat satu-ke-satu jika dan hanya jika mempertahankan 
kebebasan linear.
    \begin{answer}
      Kita membuktikan bahwa \( \map{h}{V}{W} \) bersifat satu-ke-satu jika dan hanya
      jika, untuk setiap himpunan bagian bebas linear \( S \) dari \( V \), himpunan bagian
      \( h(S)=\set{h(\vec{s})\suchthat \vec{s}\in S} \) dari \( W \) juga
      bebas linear.

      Satu arahnya mudah\Dash menurut \nearbytheorem{th:OOHomoEquivalence}, 
      jika \( h \) tidak satu-ke-satu, maka ruang nolnya
      nontrivial, yaitu memuat lebih dari sekadar vektor nol.
      Jadi, 
      jika \( \vec{v}\neq\zero_V \) berada dalam ruang nol tersebut, himpunan singleton
      \( \set{\vec{v\,}} \) bebas, sedangkan citranya
      \( \set{h(\vec{v})}=\set{\zero_W} \) tidak bebas.

      Untuk arah lainnya,
      anggaplah \( h \) satu-ke-satu sehingga menurut
      \nearbytheorem{th:OOHomoEquivalence} memiliki ruang nol trivial.
      Maka untuk sembarang \( \vec{v}_1,\dots,\vec{v}_n\in V \), relasi
      \begin{equation*}
        \zero_W=
        c_1\cdot h(\vec{v}_1)+\dots+c_n\cdot h(\vec{v}_n)
        =h(c_1\cdot \vec{v}_1+\dots+c_n\cdot \vec{v}_n)
      \end{equation*}
      mengakibatkan relasi 
      \( c_1\cdot \vec{v}_1+\dots+c_n\cdot \vec{v}_n=\zero_V \).
      Dengan demikian, jika suatu himpunan bagian dari \( V \) bebas, maka citranya dalam
      \( W \) juga bebas.  

      \textit{Catatan.}
      Pernyataannya adalah bahwa suatu pemetaan linear bersifat satu-ke-satu jika dan hanya jika
      mempertahankan kebebasan untuk \emph{semua} himpunan
      (yaitu, jika suatu himpunan bebas, maka citranya juga bebas).
      Pemetaan yang tidak satu-ke-satu mungkin saja mempertahankan sebagian himpunan bebas.
      Salah satu contohnya adalah pemetaan dari $\Re^3$ ke $\Re^2$ berikut.
      \begin{equation*}
        \colvec{x \\ y \\ z}\mapsto\colvec{x+y+z \\ 0}
      \end{equation*}
      Kebebasan linear dipertahankan untuk himpunan berikut,
      \begin{equation*}
        \set{\colvec[r]{1 \\ 0 \\ 0}}\mapsto\set{\colvec[r]{1 \\ 0}}
      \end{equation*}
      dan (dalam contoh yang agak lebih rumit) juga untuk himpunan berikut, 
      \begin{equation*}
        \set{\colvec[r]{1 \\ 0 \\ 0},\colvec[r]{0 \\ 1 \\ 0}}
        \mapsto\set{\colvec[r]{1 \\ 0}}
      \end{equation*}
      (ingat bahwa dalam suatu himpunan, unsur yang berulang tidak muncul dua kali).
      Namun, terdapat himpunan-himpunan yang kebebasannya tidak dipertahankan di bawah
      pemetaan ini, 
      \begin{equation*}
        \set{\colvec[r]{1 \\ 0 \\ 0},\colvec[r]{0 \\ 2 \\ 0}}
        \mapsto\set{\colvec[r]{1 \\ 0},\colvec[r]{2 \\ 0}}
      \end{equation*}
      sehingga tidak semua himpunan mengalami pelestarian kebebasan. 
    \end{answer}
  \item \label{exer:DimRngLessImpMapOnto}
    \nearbycorollary{cor:RankDecreases} menyatakan bahwa agar terdapat homomorfisme
    pada dari ruang vektor $V$ ke ruang vektor $W$, 
    dimensi $W$ perlu kurang dari atau sama dengan 
    dimensi $V$.
    Buktikan bahwa syarat ini juga cukup;
    gunakan \nearbytheorem{th:HomoDetActOnBasis} untuk menunjukkan bahwa
    jika dimensi $W$ 
    kurang dari atau sama dengan dimensi $V$,
    maka terdapat homomorfisme pada dari $V$ ke $W$.
    \begin{answer}
      (Kita menggunakan notasi dari \nearbytheorem{th:HomoDetActOnBasis}.)
      Tetapkan basis $\sequence{\vec{\beta}_1,\dots,\vec{\beta}_n}$ untuk $V$
      dan basis 
      $\sequence{\vec{w}_1,\ldots,\vec{w}_k}$ untuk $W$.
      Jika dimensi $k$ dari $W$ kurang dari atau sama dengan dimensi $n$
      dari $V$, 
      maka teorema tersebut memberikan pemetaan linear dari $V$ ke $W$
      yang ditentukan sebagai berikut.
      \begin{equation*}
        \vec{\beta}_1\mapsto\vec{w}_1,\,\dots,\,\vec{\beta}_k\mapsto\vec{w}_k
        \quad\text{dan}\quad
        \vec{\beta}_{k+1}\mapsto\vec{w}_k,
           \,\dots,\,\vec{\beta}_n\mapsto\vec{w}_k
      \end{equation*}
      Kita hanya perlu membuktikan bahwa pemetaan ini pada.

      Kita dapat menuliskan setiap anggota $W$ sebagai kombinasi linear 
      unsur-unsur basis, 
      $c_1\cdot \vec{w}_1+\dots+c_k\cdot \vec{w}_k$.
      Vektor ini merupakan citra, di bawah pemetaan yang dijelaskan di atas, dari 
      $c_1\cdot \vec{\beta}_1+\dots+c_k\cdot \vec{\beta}_k
      +0\cdot \vec{\beta}_{k+1}\dots+0\cdot \vec{\beta}_n$.
      Dengan demikian, pemetaan tersebut pada.
    \end{answer}  
  \recommended \item
    Ingat bahwa ruang nol merupakan himpunan bagian dari domain dan ruang hasil 
    merupakan himpunan bagian dari kodomain. 
    Apakah keduanya mesti berbeda?
    Apakah terdapat homomorfisme yang ruang nol dan ruang hasilnya
    memiliki irisan nontrivial?
    \begin{answer}
      Ya.
      Untuk transformasi pada \( \Re^2 \) yang diberikan oleh
      \begin{equation*}
        \colvec{x \\ y}\mapsunder{h}\colvec{0 \\ x}
      \end{equation*}
      kita memperoleh hasil berikut.
      \begin{equation*}
        \nullspace{h}=\set{\colvec{0 \\ y}\suchthat y\in\Re}
        =\rangespace{h}
      \end{equation*}
      \textit{Catatan.}
      Kita akan melihat lebih banyak mengenai hal ini dalam bab kelima.   
    \end{answer}
  \item 
        Buktikan bahwa 
        citra suatu rentang sama dengan rentang citra-citranya.
        Dengan kata lain, jika \( \map{h}{V}{W} \) linear, 
        buktikan bahwa jika \( S \) merupakan himpunan bagian dari 
        \( V \), maka \( h(\spanof{S}) \) sama dengan \( \spanof{h(S)} \).
        Hasil ini memperumum \nearbylemma{le:RangeIsSubSp}
        karena menunjukkan bahwa jika \( U \) adalah sembarang subruang dari \( V \), maka
        citranya \( \set{h(\vec{u})\suchthat \vec{u}\in U} \)
        merupakan subruang dari \( W \), sebab rentang himpunan $U$ adalah $U$.
    \begin{answer}
        Berikut perhitungan sederhana.
          \begin{align*}
            h(\spanof{S})
            &=\set{h(c_1\vec{s}_1+\dots+c_n\vec{s}_n)
                 \suchthat \text{\( c_1,\dots,c_n\in\Re \)
                                 dan \( \vec{s}_1,\dots,\vec{s}_n\in S \)} } \\
            &=\set{c_1h(\vec{s}_1)+\dots+c_nh(\vec{s}_n)
                 \suchthat \text{\( c_1,\dots,c_n\in\Re \)
                                 dan \( \vec{s}_1,\dots,\vec{s}_n\in S \)} } \\
            &=\spanof{h(S)}
          \end{align*}   
     \end{answer}
  \item \label{exer:Cosets}   
    \begin{exparts}
    \partsitem Buktikan bahwa untuk sembarang pemetaan linear \( \map{h}{V}{W} \) dan sembarang
      \( \vec{w}\in W \), himpunan \( h^{-1}(\vec{w}) \) berbentuk
      \begin{equation*}
        h^{-1}(\vec{w})=
        \set{\vec{v}+\vec{n}\suchthat \text{$\vec{v},\vec{n}\in V$,
                                            $\vec{n}\in\nullspace{h}$,
                                            dan $h(\vec{v})=\vec{w}$}}
      \end{equation*}
      (jika \( h \) tidak pada dan $\vec{w}$ tidak berada dalam ruang hasil $h$,
      maka himpunan ini kosong karena syarat ketiganya tidak dapat dipenuhi).
      Himpunan semacam ini merupakan \definend{koset}\index{koset} dari \( \nullspace{h} \)
      dan kita menyatakannya sebagai \( \vec{v}+\nullspace{h} \).
    \partsitem Tinjau pemetaan \( \map{t}{\Re^2}{\Re^2} \) yang diberikan oleh
      \begin{equation*}
        \colvec{x \\ y}
          \mapsunder{t}
        \colvec{ax+by \\ cx+dy}
      \end{equation*}
      untuk suatu skalar $a$, $b$, $c$, dan $d$.
      Buktikan bahwa \( t \) linear.
    \partsitem
      Simpulkan dari kedua butir sebelumnya bahwa untuk sembarang sistem linear berbentuk
      \begin{equation*}
        \begin{linsys}{2}
          ax  &+  &by  &=  &e \\
          cx  &+  &dy  &=  &f 
        \end{linsys}
      \end{equation*}
      kita dapat menuliskan himpunan solusinya
      (vektor-vektornya merupakan anggota $\Re^2$)
      \begin{equation*}
        \set{\vec{p}+\vec{h}\suchthat
             \text{\( \vec{h} \) memenuhi sistem homogen terkait} }
      \end{equation*}
      dengan \( \vec{p} \) suatu solusi partikular dari sistem linear tersebut
      (jika tidak terdapat solusi partikular, maka himpunan di atas
      kosong).
    \partsitem Tunjukkan bahwa pemetaan
      \( \map{h}{\Re^n}{\Re^m} \) berikut linear,
      \begin{equation*}
        \colvec{x_1 \\ \vdots \\ x_n}
          \mapsto
        \colvec{a_{1,1}x_1+\dots+a_{1,n}x_n \\ 
                     \vdots \\ 
                     a_{m,1}x_1+\dots+a_{m,n}x_n}
      \end{equation*}
      untuk sembarang skalar \( a_{1,1} \), \ldots, \( a_{m,n} \). 
      Perluas kesimpulan yang diperoleh pada butir sebelumnya.
    \partsitem Tunjukkan bahwa pemetaan turunan ke-\( k \) merupakan transformasi 
      linear pada
      \( \polyspace_n \) untuk setiap \( k \).
      Buktikan bahwa pemetaan berikut
      merupakan transformasi linear pada ruang tersebut,
      \begin{equation*}
        f\mapsto
        \frac{d^k}{dx^k}f+c_{k-1}\frac{d^{k-1}}{dx^{k-1}}f
           +\dots+
        c_1\frac{d}{dx}f+c_0f
      \end{equation*}
      untuk sembarang skalar \( c_{k-1} \), \ldots, \( c_0 \).
      Tarik kesimpulan seperti di atas.
    \end{exparts}
    \begin{answer}
      \begin{exparts}
        \partsitem 
          Kita akan menunjukkan bahwa himpunan-himpunan berikut sama,
          $h^{-1}(\vec{w})
            =\set{\vec{v}+\vec{n}\suchthat \vec{n}\in\nullspace{h} }$ 
          melalui inklusi timbal balik.
          Untuk arah 
          $\set{\vec{v}+\vec{n}\suchthat \vec{n}\in\nullspace{h} }
            \subseteq h^{-1}(\vec{w})$
          perhatikan saja bahwa 
          $h(\vec{v}+\vec{n})=h(\vec{v})+h(\vec{n})$ sama dengan
          $\vec{w}$, sehingga setiap anggota himpunan pertama merupakan anggota himpunan
          kedua. 
          Untuk arah 
          $h^{-1}(\vec{w})
            \subseteq\set{\vec{v}+\vec{n}\suchthat \vec{n}\in\nullspace{h} }$ 
          tinjau $\vec{u}\in h^{-1}(\vec{w})$.
          Karena \( h \) linear, \( h(\vec{u})=h(\vec{v}) \)
          mengakibatkan \( h(\vec{u}-\vec{v})=\zero \).
          Kita dapat menuliskan \( \vec{u}-\vec{v} \) sebagai \( \vec{n} \),
          lalu kita memperoleh
          $\vec{u}\in\set{\vec{v}+\vec{n}\suchthat \vec{n}\in\nullspace{h} }$,
          seperti yang diinginkan, karena $\vec{u}=\vec{v}+(\vec{u}-\vec{v})$.
        \partsitem Pemeriksaan ini rutin.
        \partsitem Hasil ini langsung.
        \partsitem Secara ringkas, untuk pemeriksaan linearitas, 
          dengan \( c,d \) skalar dan
          \( \vec{x},\vec{y}\in\Re^n \) berkomponen
          \( x_1,\dots,x_n \) dan \( y_1,\dots,y_n \),
          kita memperoleh hasil berikut.
          \begin{align*}
            h(c\cdot \vec{x}+d\cdot \vec{y})
            &=\colvec{a_{1,1}(cx_1+dy_1)+\dots+a_{1,n}(cx_n+dy_n) \\ 
                         \vdots \\ 
                         a_{m,1}(cx_1+dy_1)+\dots+a_{m,n}(cx_n+dy_n)}    \\
            &=\colvec{a_{1,1}cx_1+\dots+a_{1,n}cx_n \\ 
                         \vdots \\ 
                         a_{m,1}cx_1+\dots+a_{m,n}cx_n}
            +\colvec{a_{1,1}dy_1+\dots+a_{1,n}dy_n \\ 
                         \vdots \\                           
                         a_{m,1}dy_1+\dots+a_{m,n}dy_n}        \\
            &=c\cdot h(\vec{x})+d\cdot h(\vec{y})
          \end{align*}
          Kesimpulan yang sesuai adalah
          \( \text{Umum}=\text{Partikular}+\text{Homogen} \).
        \partsitem Setiap pangkat operator turunan bersifat 
          linear karena aturan
          \begin{equation*}
            \frac{d^k}{dx^k}(f(x)+g(x))=\frac{d^k}{dx^k}f(x)
                                         +\frac{d^k}{dx^k}g(x)
            \quad\text{dan}\quad
            \frac{d^k}{dx^k}rf(x)=r\frac{d^k}{dx^k}f(x)
          \end{equation*}
          dari kalkulus.
          Dengan demikian, pemetaan yang diberikan merupakan transformasi linear pada ruang tersebut karena
          setiap kombinasi linear pemetaan-pemetaan linear juga merupakan pemetaan linear
          menurut \nearbylemma{le:SpLinFcns}.
          Kesimpulan yang sesuai adalah 
          \( \text{Umum}=\text{Partikular}+\text{Homogen} \),
          dengan persamaan diferensial homogen terkait 
          memiliki konstanta \( 0 \).
      \end{exparts}  
     \end{answer}
  \item 
    Buktikan bahwa untuk sembarang transformasi \( \map{t}{V}{V} \) dengan rank satu, 
    pemetaan yang diperoleh dengan mengomposisikan operator tersebut dengan dirinya sendiri, 
    \( \map{\composed{t}{t}}{V}{V} \), memenuhi
    \( \composed{t}{t}=r\cdot t \) untuk suatu bilangan real \( r \).
    \begin{answer}
      Karena rank \( t \) adalah satu, ruang hasil \( t \)
      merupakan subruang satu dimensi.
      Dengan mengambil $\sequence{t(\vec{v})}$ sebagai basis (untuk suatu
      $\vec{v}$ yang sesuai), untuk setiap $\vec{w}\in V$ kita memperoleh bahwa citra 
      $t(\vec{w})\in V$ merupakan kelipatan vektor basis ini\Dash terkait
      dengan setiap $\vec{w}$ terdapat skalar 
      \( c_{\vec{w}} \)  
      sedemikian sehingga \( t(\vec{w})=c_{\vec{w}}t(\vec{v}) \).
      Terapkan \( t \) pada kedua ruas persamaan tersebut 
      dan ambil \( r \) sebagai \( c_{t(\vec{v})} \),
      \begin{equation*}
        \composed{t}{t}(\vec{w})
        =t(c_{\vec{w}}\cdot t(\vec{v}))
        =c_{\vec{w}}\cdot\composed{t}{t}(\vec{v})
        =c_{\vec{w}}\cdot c_{t(\vec{v})}\cdot t(\vec{v})
        =c_{\vec{w}}\cdot r\cdot t(\vec{v})
        =r\cdot c_{\vec{w}}\cdot t(\vec{v})
        =r\cdot t(\vec{w})
      \end{equation*}
      untuk memperoleh kesimpulan yang diinginkan.  
    \end{answer}
  \item 
    Misalkan \( \map{h}{V}{\Re} \) suatu homomorfisme, tetapi bukan
    homomorfisme nol.
    Buktikan bahwa jika \( \sequence{\vec{\beta}_1,\dots,\vec{\beta}_n} \) merupakan
    basis ruang nol dan jika \( \vec{v}\in V \) tidak berada dalam ruang nol,
    maka \( \sequence{\vec{v},\vec{\beta}_1,\dots,\vec{\beta}_n} \) merupakan
    basis seluruh domain \( V \).
    \begin{answer}
      Menurut asumsi, \( h \) bukan pemetaan nol, 
      sehingga terdapat vektor \( \vec{v}\in V \) yang tidak berada dalam ruang nol.
      Perhatikan bahwa \( \sequence{h(\vec{v})} \) merupakan basis \( \Re \) 
      karena merupakan himpunan bagian bebas linear berukuran satu dari \( \Re \).
      Akibatnya, \( h \) bersifat pada, sebab untuk sembarang \( r\in\Re \), kita memiliki
      \( r=c\cdot h(\vec{v}) \) untuk suatu skalar \( c \), sehingga
      \( r=h(c\vec{v}) \).

      Dengan demikian, rank \( h \) adalah satu.
      Karena nulitasnya adalah $n$, dimensi domain
      \( h \), yaitu ruang vektor \( V \), adalah \( n+1 \).
      Kita dapat menuntaskan bukti dengan menunjukkan bahwa
      \( \set{\vec{v},\vec{\beta}_1,\dots,\vec{\beta}_n} \)
      bebas linear, karena himpunan itu berukuran~$n+1$ dalam ruang 
      berdimensi~$n+1$.
      Karena \( \set{\vec{\beta}_1,\dots,\vec{\beta}_n} \) bebas
      linear, kita hanya perlu menunjukkan bahwa 
      \( \vec{v} \) bukan kombinasi linear vektor-vektor lainnya.
      Namun, \( c_1\vec{\beta}_1+\dots+c_n\vec{\beta}_n=\vec{v} \) akan memberikan
      \( -\vec{v}+c_1\vec{\beta}_1+\dots+c_n\vec{\beta}_n=\zero \), dan menerapkan
      \( h \) pada kedua ruas akan menghasilkan kontradiksi.  
    \end{answer}
  \item 
    Tunjukkan bahwa untuk sembarang ruang \( V \) berdimensi \( n \),
    \definend{ruang dual}\index{ruang dual}\index{ruang vektor!dual}
    \begin{equation*}
      \linmaps{V}{\Re}=\set{\map{h}{V}{\Re}\suchthat \text{\( h \) linear}}
    \end{equation*}
    isomorfik dengan \( \Re^n \).
    Ruang ini sering dinyatakan dengan $V^\ast$.
    Simpulkan bahwa \( V^\ast\isomorphicto V \).
    \begin{answer}
      Tetapkan basis \( \sequence{\vec{\beta}_1,\dots,\vec{\beta}_n} \)
      untuk \( V \).
      Kita akan membuktikan bahwa pemetaan
      \begin{equation*}
        h\mapsunder{\Phi}
        \colvec{h(\vec{\beta}_1) \\ \vdots \\ h(\vec{\beta}_n)}
      \end{equation*}
      merupakan isomorfisme dari \( V^\ast \) ke \( \Re^n \).

      Untuk melihat bahwa $\Phi$ satu-ke-satu, anggaplah $h_1$ dan $h_2$ merupakan
      anggota $V^\ast$ sedemikian sehingga $\Phi(h_1)=\Phi(h_2)$.
      Maka
      \begin{equation*}
        \colvec{h_1(\vec{\beta}_1) \\ \vdots \\ h_1(\vec{\beta}_n)}
        =\colvec{h_2(\vec{\beta}_1) \\ \vdots \\ h_2(\vec{\beta}_n)}
      \end{equation*}
      dan akibatnya, $h_1(\vec{\beta}_1)=h_2(\vec{\beta}_1)$, dan seterusnya.
      Namun, suatu homomorfisme ditentukan oleh aksinya pada basis, sehingga 
      \( h_1=h_2 \), dan oleh karena itu \( \Phi \) satu-ke-satu.

      Untuk melihat bahwa $\Phi$ pada, tinjau
      \begin{equation*}
        \colvec{x_1 \\ \vdots \\ x_n}
      \end{equation*}
      untuk \( x_1,\ldots,x_n\in\Re \). 
      Fungsi $h$ dari $V$ ke $\Re$ berikut 
      \begin{equation*}
        c_1\vec{\beta}_1+\dots+c_n\vec{\beta}_n
        \mapsunder{h} c_1x_1+\dots+c_nx_n
      \end{equation*}
      bersifat linear dan $\Phi$ memetakannya ke 
      vektor yang diberikan dalam $\Re^n$, sehingga \( \Phi \) pada.

      Pemetaan \( \Phi \) juga mempertahankan struktur:~dengan
      \begin{align*}
        c_1\vec{\beta}_1+\dots+c_n\vec{\beta}_n
        &\mapsunder{h_1}
        c_1h_1(\vec{\beta}_1)+\dots+c_nh_1(\vec{\beta}_n)    \\
        c_1\vec{\beta}_1+\dots+c_n\vec{\beta}_n
        &\mapsunder{h_2}
        c_1h_2(\vec{\beta}_1)+\dots+c_nh_2(\vec{\beta}_n)
      \end{align*}
      kita memperoleh
      \begin{multline*}
        (r_1h_1+r_2h_2)(c_1\vec{\beta}_1+\dots+c_n\vec{\beta}_n)       \\
        \begin{aligned}
        &=
        c_1(r_1h_1(\vec{\beta}_1)+r_2h_2(\vec{\beta}_1))
          +\dots
          +c_n(r_1h_1(\vec{\beta}_n)+r_2h_2(\vec{\beta}_n))    \\
        &=
        r_1(c_1h_1(\vec{\beta}_1)+\dots+c_nh_1(\vec{\beta}_n))
          + r_2(c_1h_2(\vec{\beta}_1)+\dots+c_nh_2(\vec{\beta}_n))
        \end{aligned}
      \end{multline*}
      sehingga \( \Phi(r_1h_1+r_2h_2)=r_1\Phi(h_1)+r_2\Phi(h_2) \).  
    \end{answer}
  \item 
    Tunjukkan bahwa setiap pemetaan linear merupakan jumlah pemetaan-pemetaan ber-rank satu.
    \begin{answer}
      Misalkan \( \map{h}{V}{W} \) linear dan tetapkan basis
      \( \sequence{\vec{\beta}_1,\dots,\vec{\beta}_n} \) untuk \( V \).
      Tinjau \( n \) pemetaan berikut dari \( V \) ke \( W \),
      \begin{equation*}
        h_1(\vec{v})=c_1\cdot h(\vec{\beta}_1),
        \quad
        h_2(\vec{v})=c_2\cdot h(\vec{\beta}_2),
        \quad\ldots\quad,
        h_n(\vec{v})=c_n\cdot h(\vec{\beta}_n)
      \end{equation*}
      untuk sembarang \( \vec{v}=c_1\vec{\beta}_1+\dots+c_n\vec{\beta}_n \).
      Jelas bahwa \( h \) merupakan jumlah dari pemetaan-pemetaan \( h_i \).
      Kita hanya perlu memeriksa bahwa setiap \( h_i \) linear:~jika
      \( \vec{u}=d_1\vec{\beta}_1+\dots+d_n\vec{\beta}_n \), maka
      \( h_i(r\vec{v}+s\vec{u})=(rc_i+sd_i)h(\vec{\beta}_i)
      =rh_i(\vec{v})+sh_i(\vec{u}) \).  
    \end{answer}
 \item 
   Apakah `homomorfik ke' merupakan relasi ekuivalensi?
   (\textit{Petunjuk:}~kesulitannya adalah menentukan makna yang tepat 
   untuk frasa dalam tanda kutip tersebut.)
   \begin{answer}
     Bisa ya (secara trivial) atau tidak (hampir secara trivial).

     Jika kita memaknai \( V \) `homomorfik ke' \( W \) sebagai adanya
     homomorfisme dari \( V \) ke dalam (tetapi tidak harus pada terhadap) \( W \),
     maka setiap ruang homomorfik ke setiap ruang lainnya karena pemetaan nol selalu
     ada.

     Jika kita memaknai \( V \) `homomorfik ke' \( W \) sebagai adanya
     homomorfisme pada dari \( V \) ke \( W \), maka relasi tersebut bukan
     relasi ekuivalensi.
     Sebagai contoh, terdapat homomorfisme
     pada dari \( \Re^3 \) ke \( \Re^2 \)
     (proyeksi adalah salah satunya), tetapi menurut \nearbycorollary{cor:RankDecreases}, tidak ada homomorfisme dari \( \Re^2 \) yang pada terhadap
     \( \Re^3 \),
     sehingga relasi tersebut tidak simetris.\appendrefs{relasi ekuivalensi}  
   \end{answer}
  \item 
    Tunjukkan bahwa ruang hasil dan ruang nol dari pangkat-pangkat pemetaan linear
    \( \map{t}{V}{V} \) masing-masing membentuk rantai menurun,
    \begin{equation*}
      V\supseteq \rangespace{t}\supseteq\rangespace{t^2}\supseteq\ldots
    \end{equation*}
    dan rantai menaik,
    \begin{equation*}
     \set{\vec{0}}\subseteq\nullspace{t}\subseteq\nullspace{t^2}\subseteq\ldots
    \end{equation*}
    Tunjukkan pula bahwa jika \( k \) sedemikian sehingga
    \( \rangespace{t^k}=\rangespace{t^{k+1}} \)
    maka semua ruang hasil berikutnya sama:
    \( \rangespace{t^k}=\rangespace{t^{k+1}}=\rangespace{t^{k+2}}\ldots\, \).
    Serupa dengan itu, jika \( \nullspace{t^k}=\nullspace{t^{k+1}} \),
    maka  
    \( \nullspace{t^k}=\nullspace{t^{k+1}}=\nullspace{t^{k+2}}=\ldots\, \).
    \begin{answer}
      Fakta bahwa ruang-ruang tersebut membentuk rantai sudah jelas.
      Untuk bagian lainnya, di sini kita menunjukkan bahwa
      \( \rangespace{t^{j+1}}=\rangespace{t^{j}} \)
      mengakibatkan
      \( \rangespace{t^{j+2}}=\rangespace{t^{j+1}} \).
      Selanjutnya, induksi dapat diterapkan.

      Anggaplah 
      \( \rangespace{t^{j+1}}=\rangespace{t^{j}} \).
      Maka \( \map{t}{\rangespace{t^{j+1}}}{\rangespace{t^{j+2}}} \)
      adalah pemetaan yang sama, dengan domain yang sama, seperti
      \( \map{t}{\rangespace{t^{j}}}{\rangespace{t^{j+1}}} \).
      Dengan demikian, pemetaan itu memiliki ruang hasil yang sama:
      \( \rangespace{t^{j+2}}=\rangespace{t^{j+1}} \).  
    \end{answer}
\end{exercises}
